\documentclass[11pt,a4paper]{ctexart}
\usepackage[margin=2.2cm]{geometry}
\usepackage{amsmath,amssymb,booktabs,longtable,graphicx,float,xcolor,adjustbox,hyperref}
\usepackage[section]{placeins}
\hypersetup{colorlinks=true,linkcolor=blue!50!black,citecolor=blue!50!black,urlcolor=blue!50!black}
\setlength{\LTcapwidth}{\textwidth}
\newcommand{\Mtwo}{方法二}
\title{方法二（FD-LGF[FDM 背景]）与传统 LGF、GMG、BoomerAMG 的\\二维和三维对比}
\author{温州肯恩大学 \quad Nonlinear Poisson 项目}
\date{2026 年 11 月 3 日}

\begin{document}
\maketitle

\begin{abstract}
对 $-\nabla\!\cdot(a\nabla u)+qu=f$ 的通量格式离散，比较四种 PCG 预条件：传统 LGF（常系数格点 Green 函数）、几何多重网格 GMG、
hypre BoomerAMG，以及\Mtwo{}——分层背景 $a_1(x)$ 的精确快速逆加夹杂附近的精确局部 LU，按平衡形式组合。
二维取 $n=127$–$1023$（最多 $10^6$ 个未知量），三维取 $n=31$–$127$（最多 $2\times10^6$ 个未知量），对比度
$\chi=10^2$–$10^8$，四类系数。统计迭代次数、精度（代数误差、相对连续解的总误差、逐步误差曲线）和时间（setup、solve、
达到误差 $10^{-8}$ 的时间）。
\textbf{主要结论：}(1) 在教授的 $a=\chi^x$ 和阶跃分层上，\Mtwo{} 在二维、三维都只需 1--3 次迭代，与 $\chi$、$n$ 无关，
最大网格上的总时间二维不到 BoomerAMG 的 0.3 倍，三维不到 0.06 倍；传统 LGF 在 $\chi\ge10^4$ 时不收敛。
(2) 有夹杂时，\Mtwo{} 的迭代数仍与 $\chi$、$n$ 无关（二维 7--10 次，三维 10--13 次），时间则取决于夹杂区稀疏 LU 的规模：
二维 $n=1023$ 时比 BoomerAMG 慢 1.4--3.4 倍；三维 $n=127$ 时，规则夹杂上两者持平（0.92--1.03 倍），
随机夹杂上\Mtwo{} 快约 2 倍（0.48--0.52 倍）。三维 BoomerAMG 的迭代数随 $n$ 增长（11$\to$17），\Mtwo{} 不增长。
(3) 所有收敛方法的离散误差相同；按残差停止时，\Mtwo{} 在低系数区的代数误差比 BoomerAMG 大 1--4 个数量级，
需要更紧的容差或按误差计数。
\end{abstract}

\tableofcontents

\section{问题与离散}
在 $\Omega=(0,1)^d$（$d=2,3$）上求解
\[
  -\nabla\cdot\big(a(\mathbf x)\nabla u\big)+q(\mathbf x)\,u=f,\qquad u|_{\partial\Omega}=0 .
\]
网格 $h=1/(n+1)$，每个方向 $n$ 个内点。采用守恒的 $2d+1$ 点通量格式，面系数取两端节点的算术平均
$a_{i+1/2}=\tfrac12(a_i+a_{i+1})$（与 Xia 论文式 (5)--(6) 相同；三维同样按三个方向写出），得到对称正定系统 $Au=f$。
所有方法共用同一个 PCG 循环（\texttt{axy\_lgf.solvers.\_pcg} 及其三维复制）：
停止规则为 $\|r_k\|/\|f\|\le10^{-10}$（收敛时重算真实残差），二维 maxiter $=1000$，三维 $500$。
计时为单线程（\texttt{OMP\_NUM\_THREADS=1}），同一台 8 核 Mac，取 3 次中位数（单次 $>5$ s 时只跑 1 次）。

\section{四种方法}
\subsection{传统 LGF}
按 Xia 论文 §4.1 式 (14)--(15)：$\bar a=\mathrm{mean}(a)$，$\lambda=\mathrm{mean}(\max(\sigma_h+q_u,0))$。
本文的问题是线性的（$a$ 不依赖 $u$，$\sigma_h=0$，$q_u=q$），所以 $\lambda=\mathrm{mean}(q)$。预条件为
\[
  M_{\mathrm{LGF}}^{-1}=\big(\bar a(-\Delta_h)+\lambda I\big)^{-1},
\]
在方形区域上用 $d$ 维 DST-I 精确作用（论文中 BAE 处理不规则边界；在方形区域上边界与辅助盒子重合，BAE 部分退化为恒等）。
论文的 Jacobian 非对称，外层用 FGMRES；本文矩阵对称正定，外层用 PCG。数值检查：常系数时与真逆的相对误差 $10^{-14}$。

\subsection{GMG}
双线性（三维为三线性）插值 $P$，限制 $P^{\mathsf T}$，Galerkin 粗算子 $P^{\mathsf T}AP$，V(2,2) 阻尼 Jacobi，
最粗层（$3^d$ 个点）Cholesky 直接解。阻尼取 $\omega_\ell=c/\max_i\sum_j|D^{-1}A|_{ij}$，
二维 $c=1.6$、三维 $c=12/7$，在 Laplace 上正好是教科书的最优值 $4/5$（二维）、$6/7$（三维）
（Briggs--Henson--McCormick 2000；Trottenberg 等 2001）。由 Gershgorin，$\omega_\ell\lambda_{\max}(D^{-1}A)\le c<2$，
所以光滑子收敛、V-cycle 对称正定（数值检查对称到 $10^{-18}$）。
\emph{说明：}项目原 GMG 取 $c=1$（Laplace 上 $\omega=1/2$），二维 Poisson 上 V(2,2) 收敛因子 0.30，而 $\omega=4/5$ 约为 0.13；
经确认后改为教科书值并重跑了二维 GMG 的全部 64 组（旧结果保存在 \texttt{results/benchmark\_gmg\_omega05.csv}）。

\subsection{BoomerAMG}
hypre BoomerAMG（Henson \& Yang 2002）经 petsc4py 调用，每次作用一个 V-cycle。参数：Falgout 粗化，classical 插值，
symmetric-SOR/Jacobi 光滑（hypre 文档要求作 CG 预条件时用对称光滑子），粗网格 Gaussian elimination，
强连接阈值二维 0.25、三维 0.5（hypre 对二维、三维的推荐值）。数值检查：对称到 $10^{-17}$，二维 Poisson 上每个 V-cycle 收敛因子 0.078。

\subsection{\Mtwo{}：FD-LGF[FDM 背景]}
把系数写成 $a=b\cdot\phi$，其中 $b=a_1(x)$ 是分层背景（本文所有算例的背景只沿 $x$ 变化，含 $\chi^x$ 和阶跃），
$\phi$ 只在夹杂内不等于 1。
\begin{itemize}
\item \textbf{背景逆 $G$}：用系数 $b$ 组装的通量算子 $B$ 可以精确快速求逆：在 $y$（三维为 $y,z$）方向做 DST-I，
  每个横向模态在 $x$ 方向解一个三对角方程，复杂度 $O(n^d\log n)$。$q$ 用其横向平均代替，是 $G$ 中唯一的近似。
  对一般可分离 $a_1(x)a_2(y)$ 用快速对角化（Lynch--Rice--Thomas 1964）。
\item \textbf{局部解 $Q$}：特征集 $\{|\log a-\log b|>\log 2\}$ 向外膨胀宽度 $w$ 的一圈，对每个连通区域取 $A$ 的子块做稀疏 LU，
  $Q=\sum_k R_k^{\mathsf T}A_k^{-1}R_k$。
\item \textbf{平衡组合}：
  \[ P=Q+(I-QA)\,G\,(I-AQ). \]
  $G$、$A_k$ 对称正定，所以 $P$ 对称正定；没有夹杂时 $P=G$。
\end{itemize}
环宽：二维 $w=1/16$（主设置，与 H30 相同）及 $1/32$、$1/64$；三维 $w=1/32$、$1/64$（$w=1/16$ 时 27 个球连成一片，
$n=127$ 的三维 LU 不可行）。

\section{算例与指标}
\begin{itemize}
\item[A] $a=\chi^x$（教授指定的对比度算例），$u=\prod\sin(\pi x_i)$，$q=1+0.2\prod\cos(2\pi x_i)$，$f$ 由连续解析式给出，
  因此可以计算相对 PDE 解的总误差。$u_h$ 二维用稀疏直接法，三维用\Mtwo{}-PCG 迭代到舍入误差（与稀疏直接法相比：$n=31$ 时差 $2\times10^{-15}$，$n=63$ 时差 $2\times10^{-14}$）。
\item[B] 阶跃层 $a=\chi^{H(x-1/2)}$，离散 MMS（$f=A_hu^*$，$u_h=u^*$）。
\item[C] $\chi^x$ 乘规则夹杂（二维 9 个圆、三维 27 个球，$r=0.08$，系数再乘 $10^4$），离散 MMS。
\item[D] $\chi^x$ 乘随机夹杂（二维 20 个圆、三维 30 个球，系数乘 $10^4$ 或 $10^{-4}$，即导体和绝缘体都有），离散 MMS。
\end{itemize}
指标：迭代次数；按残差规则停止时的代数误差 $\|u-u_h\|/\|u_h\|$；相对连续解的总误差（A）；
同一 PCG 递推逐步记录的误差曲线，由此得到第 5 步误差、误差降到 $10^{-8}$ 所需迭代 $k_8$、
以及“达到误差 $10^{-8}$ 的时间”$=$ setup $+k_8\times$单步时间；setup、solve 和总时间。
实验方案和预测在运行前锁定（SHA-256：二维 \texttt{38eab086}，三维 \texttt{06cc1590}）。

\section{二维结果}
\begin{table}[H]\centering\small
\caption{二维：PCG 迭代次数（$\chi\in\{10^2,10^4,10^6,10^8\}$ 的最小--最大值；$^{\times k}$ 表示有 $k$ 个 $\chi$ 在 1000 次内未收敛）}\label{tab:its2d}
\begin{adjustbox}{max width=\textwidth}
\begin{tabular}{llcccccc}
\toprule
算例 & $n$ & 传统 LGF & GMG & BoomerAMG & M2 $w{=}\frac1{16}$ & M2 $w{=}\frac1{32}$ & M2 $w{=}\frac1{64}$ \\
\midrule
A: $\chi^x$ & 127 & 110--1000$^{\times2}$ & 7 & 7 & 2--3 &  &  \\
 & 255 & 117--1000$^{\times2}$ & 7--8 & 7 & 2--3 &  &  \\
 & 511 & 123--1000$^{\times2}$ & 7--8 & 7 & 2--3 &  &  \\
 & 1023 & 127--1000$^{\times3}$ & 7--8 & 7--8 & 2--3 &  &  \\
\midrule
B: 阶跃层 & 127 & 15--51 & 7 & 8 & 1--3 &  &  \\
 & 255 & 15--53 & 7 & 8 & 1--3 &  &  \\
 & 511 & 17--55 & 7 & 8 & 1--3 &  &  \\
 & 1023 & 19--58 & 7 & 9 & 1--3 &  &  \\
\midrule
C: $\chi^x$+规则夹杂 & 127 & 1000$^{\times4}$ & 48--58 & 9--10 & 8 & 10--11 & 14--15 \\
 & 255 & 1000$^{\times4}$ & 56--68 & 9--10 & 8 & 10--11 & 14--15 \\
 & 511 & 1000$^{\times4}$ & 82--87 & 10 & 8 & 10--11 & 13--14 \\
 & 1023 & 1000$^{\times4}$ & 90--95 & 10--11 & 7--8 & 10--11 & 13--14 \\
\midrule
D: $\chi^x$+随机夹杂 & 127 & 1000$^{\times4}$ & 73--80 & 8--9 & 9--10 & 13--15 & 19--21 \\
 & 255 & 1000$^{\times4}$ & 87--93 & 10 & 9--10 & 13--14 & 19--21 \\
 & 511 & 1000$^{\times4}$ & 98--104 & 10 & 9--10 & 13--14 & 19--20 \\
 & 1023 & 1000$^{\times4}$ & 108--114 & 10--11 & 8--10 & 13--14 & 19--20 \\
\bottomrule
\end{tabular}
\end{adjustbox}
\end{table}

\begin{table}[H]\centering\small
\caption{二维：误差 $\|u_k-u_h\|/\|u_h\|\le10^{-8}$ 所需迭代数（各 $\chi$ 的范围；--：1000 次内未达到）}\label{tab:k82d}
\begin{adjustbox}{max width=\textwidth}
\begin{tabular}{llcccccc}
\toprule
算例 & $n$ & 传统 LGF & GMG & BoomerAMG & M2 $w{=}\frac1{16}$ & M2 $w{=}\frac1{32}$ & M2 $w{=}\frac1{64}$ \\
\midrule
A: $\chi^x$ & 127 & 68 & 5--7 & 5 & 2--3 &  &  \\
 & 255 & 71 & 5--7 & 5 & 2--3 &  &  \\
 & 511 & 73 & 5--6 & 5 & 2--3 &  &  \\
 & 1023 & 73 & 5--6 & 5 & 2--3 &  &  \\
\midrule
B: 阶跃层 & 127 & 11--69 & 5 & 6 & 3 &  &  \\
 & 255 & 11--71 & 5 & 5--6 & 3 &  &  \\
 & 511 & 10--72 & 5 & 5--6 & 3 &  &  \\
 & 1023 & 10--72 & 4--5 & 5--6 & 3 &  &  \\
\midrule
C: $\chi^x$+规则夹杂 & 127 & -- & 55--56 & 8--9 & 9--10 & 12--13 & 15--18 \\
 & 255 & -- & 64--65 & 8--9 & 9--10 & 12--13 & 15--17 \\
 & 511 & -- & 81--84 & 9 & 9--10 & 12--13 & 15--16 \\
 & 1023 & -- & 88--89 & 9 & 9--10 & 12--13 & 15--16 \\
\midrule
D: $\chi^x$+随机夹杂 & 127 & -- & 75--77 & 7--8 & 10--11 & 16--18 & 23--26 \\
 & 255 & -- & 88--89 & 8--9 & 11 & 16--18 & 22--26 \\
 & 511 & -- & 98--100 & 9--11 & 11 & 15--18 & 22--25 \\
 & 1023 & -- & 107--108 & 9 & 11 & 15--18 & 22--25 \\
\bottomrule
\end{tabular}
\end{adjustbox}
\end{table}

\begin{table}[H]\centering\small
\caption{二维：按残差规则 $\|r\|/\|f\|\le10^{-10}$ 停止时的代数误差 $\|u-u_h\|/\|u_h\|$（各 $\chi$ 中的最大值）}\label{tab:err2d}
\begin{adjustbox}{max width=\textwidth}
\begin{tabular}{llcccccc}
\toprule
算例 & $n$ & 传统 LGF & GMG & BoomerAMG & M2 $w{=}\frac1{16}$ & M2 $w{=}\frac1{32}$ & M2 $w{=}\frac1{64}$ \\
\midrule
A: $\chi^x$ & 127 & $2.8{\times}10^{-2}$ & $3.2{\times}10^{-10}$ & $3.5{\times}10^{-12}$ & $1.4{\times}10^{-9}$ &  &  \\
 & 255 & $3.7{\times}10^{-2}$ & $3.0{\times}10^{-12}$ & $5.1{\times}10^{-12}$ & $1.4{\times}10^{-9}$ &  &  \\
 & 511 & $4.2{\times}10^{-2}$ & $1.9{\times}10^{-12}$ & $5.9{\times}10^{-12}$ & $1.4{\times}10^{-9}$ &  &  \\
 & 1023 & $4.5{\times}10^{-2}$ & $1.3{\times}10^{-12}$ & $1.1{\times}10^{-11}$ & $1.4{\times}10^{-9}$ &  &  \\
\midrule
B: 阶跃层 & 127 & $6.0{\times}10^{-7}$ & $3.3{\times}10^{-12}$ & $6.5{\times}10^{-12}$ & $3.7{\times}10^{-4}$ &  &  \\
 & 255 & $6.0{\times}10^{-7}$ & $3.0{\times}10^{-12}$ & $6.7{\times}10^{-13}$ & $3.8{\times}10^{-4}$ &  &  \\
 & 511 & $6.0{\times}10^{-7}$ & $2.6{\times}10^{-12}$ & $2.9{\times}10^{-12}$ & $3.8{\times}10^{-4}$ &  &  \\
 & 1023 & $6.0{\times}10^{-7}$ & $2.4{\times}10^{-12}$ & $1.1{\times}10^{-13}$ & $3.8{\times}10^{-4}$ &  &  \\
\midrule
C: $\chi^x$+规则夹杂 & 127 & $1.9{\times}10^{-1}$ & $1.2{\times}10^{-6}$ & $2.6{\times}10^{-9}$ & $2.5{\times}10^{-7}$ & $6.2{\times}10^{-7}$ & $4.0{\times}10^{-7}$ \\
 & 255 & $2.2{\times}10^{-1}$ & $2.7{\times}10^{-7}$ & $1.8{\times}10^{-9}$ & $2.7{\times}10^{-7}$ & $5.1{\times}10^{-7}$ & $2.2{\times}10^{-7}$ \\
 & 511 & $2.3{\times}10^{-1}$ & $9.9{\times}10^{-9}$ & $2.6{\times}10^{-10}$ & $2.3{\times}10^{-7}$ & $4.1{\times}10^{-7}$ & $7.9{\times}10^{-7}$ \\
 & 1023 & $2.3{\times}10^{-1}$ & $2.9{\times}10^{-9}$ & $9.2{\times}10^{-10}$ & $7.1{\times}10^{-6}$ & $3.8{\times}10^{-7}$ & $6.9{\times}10^{-7}$ \\
\midrule
D: $\chi^x$+随机夹杂 & 127 & $3.4{\times}10^{-1}$ & $3.6{\times}10^{-8}$ & $7.3{\times}10^{-10}$ & $7.6{\times}10^{-7}$ & $2.4{\times}10^{-6}$ & $1.8{\times}10^{-6}$ \\
 & 255 & $3.7{\times}10^{-1}$ & $1.9{\times}10^{-8}$ & $1.5{\times}10^{-9}$ & $7.7{\times}10^{-7}$ & $1.7{\times}10^{-6}$ & $1.5{\times}10^{-6}$ \\
 & 511 & $3.9{\times}10^{-1}$ & $8.2{\times}10^{-9}$ & $1.5{\times}10^{-8}$ & $6.0{\times}10^{-7}$ & $2.0{\times}10^{-6}$ & $1.3{\times}10^{-6}$ \\
 & 1023 & $4.0{\times}10^{-1}$ & $4.9{\times}10^{-9}$ & $2.5{\times}10^{-10}$ & $3.0{\times}10^{-6}$ & $1.9{\times}10^{-6}$ & $1.2{\times}10^{-6}$ \\
\bottomrule
\end{tabular}
\end{adjustbox}
\end{table}

\begin{table}[H]\centering\small
\caption{二维：总时间 setup+solve（$\chi=10^8$，单线程，ms）}\label{tab:time2d}
\begin{adjustbox}{max width=\textwidth}
\begin{tabular}{llcccccc}
\toprule
算例 & $n$ & 传统 LGF & GMG & BoomerAMG & M2 $w{=}\frac1{16}$ & M2 $w{=}\frac1{32}$ & M2 $w{=}\frac1{64}$ \\
\midrule
A: $\chi^x$ & 127 & 542 & 16 & 18 & 4 &  &  \\
 & 255 & 2.36\,s & 55 & 77 & 15 &  &  \\
 & 511 & 11.07\,s & 219 & 306 & 62 &  &  \\
 & 1023 & 46.30\,s & 899 & 1.14\,s & 251 &  &  \\
\midrule
B: 阶跃层 & 127 & 26 & 15 & 21 & 3 &  &  \\
 & 255 & 114 & 49 & 83 & 12 &  &  \\
 & 511 & 536 & 198 & 323 & 50 &  &  \\
 & 1023 & 3.04\,s & 941 & 1.35\,s & 195 &  &  \\
\midrule
C: $\chi^x$+规则夹杂 & 127 & 532 & 61 & 27 & 31 & 27 & 29 \\
 & 255 & 2.11\,s & 250 & 103 & 150 & 108 & 110 \\
 & 511 & 8.09\,s & 1.37\,s & 376 & 806 & 482 & 416 \\
 & 1023 & 38.69\,s & 8.20\,s & 1.53\,s & 4.70\,s & 2.59\,s & 2.26\,s \\
\midrule
D: $\chi^x$+随机夹杂 & 127 & 627 & 87 & 27 & 36 & 35 & 40 \\
 & 255 & 2.22\,s & 391 & 109 & 160 & 135 & 143 \\
 & 511 & 8.01\,s & 1.80\,s & 444 & 891 & 640 & 597 \\
 & 1023 & 37.92\,s & 9.33\,s & 1.62\,s & 5.02\,s & 3.15\,s & 2.72\,s \\
\bottomrule
\end{tabular}
\end{adjustbox}
\end{table}

\begin{table}[H]\centering\small
\caption{二维：误差降到 $10^{-8}$ 的时间 setup$+k_8\times$单步时间（$\chi=10^8$，ms）}\label{tab:t82d}
\begin{adjustbox}{max width=\textwidth}
\begin{tabular}{llcccccc}
\toprule
算例 & $n$ & 传统 LGF & GMG & BoomerAMG & M2 $w{=}\frac1{16}$ & M2 $w{=}\frac1{32}$ & M2 $w{=}\frac1{64}$ \\
\midrule
A: $\chi^x$ & 127 & -- & 16 & 15 & 4 &  &  \\
 & 255 & -- & 50 & 65 & 15 &  &  \\
 & 511 & -- & 184 & 260 & 62 &  &  \\
 & 1023 & -- & 752 & 974 & 251 &  &  \\
\midrule
B: 阶跃层 & 127 & 40 & 12 & 18 & 7 &  &  \\
 & 255 & 171 & 41 & 66 & 25 &  &  \\
 & 511 & 785 & 163 & 258 & 102 &  &  \\
 & 1023 & 4.20\,s & 766 & 1.01\,s & 409 &  &  \\
\midrule
C: $\chi^x$+规则夹杂 & 127 & -- & 70 & 27 & 35 & 32 & 34 \\
 & 255 & -- & 283 & 103 & 168 & 127 & 126 \\
 & 511 & -- & 1.37\,s & 355 & 882 & 558 & 476 \\
 & 1023 & -- & 8.12\,s & 1.44\,s & 5.24\,s & 2.97\,s & 2.56\,s \\
\midrule
D: $\chi^x$+随机夹杂 & 127 & -- & 89 & 25 & 41 & 43 & 51 \\
 & 255 & -- & 395 & 103 & 178 & 168 & 183 \\
 & 511 & -- & 1.80\,s & 469 & 982 & 791 & 733 \\
 & 1023 & -- & 9.24\,s & 1.53\,s & 5.64\,s & 3.82\,s & 3.31\,s \\
\bottomrule
\end{tabular}
\end{adjustbox}
\end{table}

\begin{table}[H]\centering\small
\caption{二维，算例 A（$n=1023$）：相对连续解的总误差 $\|u-u^*\|/\|u^*\|$}\label{tab:aerr二维}
\begin{tabular}{lcccc}
\toprule
$\chi$ & 传统 LGF & GMG & BoomerAMG & M2 $w{=}\frac1{16}$ \\
\midrule
$10^{2}$ & $1.7{\times}10^{-6}$ & $1.7{\times}10^{-6}$ & $1.7{\times}10^{-6}$ & $1.7{\times}10^{-6}$ \\
$10^{4}$ & $1.0{\times}10^{-5}$ & $1.0{\times}10^{-5}$ & $1.0{\times}10^{-5}$ & $1.0{\times}10^{-5}$ \\
$10^{6}$ & $2.1{\times}10^{-3}$ & $2.5{\times}10^{-5}$ & $2.5{\times}10^{-5}$ & $2.5{\times}10^{-5}$ \\
$10^{8}$ & $4.5{\times}10^{-2}$ & $4.7{\times}10^{-5}$ & $4.7{\times}10^{-5}$ & $4.7{\times}10^{-5}$ \\
\bottomrule
\end{tabular}
\end{table}

\begin{figure}[H]\centering
\includegraphics[width=\textwidth]{fig2d_iters_vs_n.png}
\caption{二维：$\chi=10^8$ 时迭代次数随 $n$ 的变化。}\label{fig:2dits}
\end{figure}
\begin{figure}[H]\centering
\includegraphics[width=\textwidth]{fig2d_error_traces.png}
\caption{二维：$n=1023$、$\chi=10^8$ 时误差 $\|u_k-u_h\|/\|u_h\|$ 随迭代的下降（前 60 步）。}\label{fig:2dtrace}
\end{figure}
\begin{figure}[H]\centering
\includegraphics[width=\textwidth]{fig2d_time_to_1e-8_vs_n.png}
\caption{二维：达到误差 $10^{-8}$ 的时间（setup $+k_8\times$单步时间，$\chi=10^8$）。}\label{fig:2dtime}
\end{figure}

\subsection{二维分析}
\paragraph{迭代次数。}在 A、B 上，\Mtwo{} 只需 1--3 次迭代，与 $\chi$、$n$ 无关，因为背景逆 $G$ 对分层系数是精确的，
剩下的只有 $q$ 的横向变化；GMG 7--8 次，BoomerAMG 7--9 次；传统 LGF 在 A 上 $\chi=10^2$ 时已要 110--127 次，
$\chi\ge10^4$（$n=1023$）或 $\chi\ge10^6$ 时 1000 次内不收敛，在 B 上能收敛但要 15--58 次。
在 C、D 上，\Mtwo{}（$w=1/16$）7--10 次，BoomerAMG 8--11 次，基本相同；GMG 48--114 次，且随 $n$ 增长；传统 LGF 全部不收敛。
\Mtwo{} 的迭代数从 $n=127$ 到 $1023$ 最多变化 2 次，与网格无关。

\paragraph{精度。}A 上所有收敛方法的总误差前 3 位有效数字完全相同（如 $n=1023$、$\chi=10^8$ 时都是 $4.7\times10^{-5}$），
说明误差由离散决定；传统 LGF 在 $\chi\ge10^6$ 时没收敛，总误差比其他方法大 1.1--950 倍（网格越细、$\chi$ 越大差距越大），
见表~\ref{tab:aerr二维}。
按残差规则停止时，\Mtwo{} 的代数误差（C、D 上 $10^{-8}$--$7\times10^{-6}$，B 上 $\chi=10^8$ 时 $3.8\times10^{-4}$）
比 BoomerAMG（$10^{-10}$ 量级）大得多：\Mtwo{} 的残差在高系数区很快变小，而低系数区的误差在残差里几乎看不到（H20 已发现）。
按误差计数（表~\ref{tab:k82d}），把误差降到 $10^{-8}$ 时，\Mtwo{} 在 A、B 上要 2--3 次（BoomerAMG 5--6 次），
在 C、D 上要 9--11 次（BoomerAMG 7--11 次）。图~\ref{fig:2dtrace} 显示，有夹杂时 BoomerAMG 每步的误差下降比\Mtwo{} 略快。

\paragraph{时间。}A、B：\Mtwo{} setup 只是一次三对角分解，$n=1023$ 时总时间 0.2--0.36 s，BoomerAMG 1.1--1.4 s，
GMG 0.86--1.0 s，传统 LGF 1--49 s；\Mtwo{} 是 BoomerAMG 的 $\le$0.29 倍、GMG 的 $\le$0.41 倍。
C、D：\Mtwo{} 的成本由夹杂区 LU 决定。$w=1/16$ 时扩张后的 9 个圆连成一片，占网格的 48\%（D 为 57\%），
$n=1023$ 时 LU setup 3.2--3.5 s，总时间 4.7--5.3 s，是 BoomerAMG（1.5--1.6 s）的 3.1--3.4 倍；
$w=1/64$ 时 2.3--2.8 s，仍是 1.4--1.8 倍。$n\le255$ 时两者相当。GMG 6.9--11.7 s，传统 LGF 约 40 s 且不收敛。

\paragraph{预测检验（二维，锁定方案 \texttt{38eab086}）。}
\begin{enumerate}
\item A、B：\Mtwo{} $\le3$ 次；BoomerAMG、GMG 6--12 次；LGF 在 A 上 $\chi\ge10^6$ 不收敛——\textbf{成立}。
\item A：收敛方法总误差 3 位相同——\textbf{成立}；LGF 在 $\chi\ge10^6$ 时误差大 $\ge10$ 倍——\textbf{部分成立}
  （$n\ge511$ 成立；$n=127$、$\chi=10^6$ 只大 1.14 倍，$n=255$、$\chi=10^6$ 大 3.7 倍）。
\item C、D：\Mtwo{} $\le11$ 次且随 $n$ 变化 $\le2$；BoomerAMG $\le12$ 次；GMG $\ge40$；LGF 不收敛——\textbf{成立}。
\item A、B，$n\ge511$：\Mtwo{} 总时间 $\le$ BoomerAMG 的 1/3——\textbf{成立}（$\le0.29$）；$\le$ GMG 的 1/3——\textbf{不成立}
  （最大 0.41，GMG 改用 $\omega=4/5$ 后更快；该预测是针对原 GMG 写的）。
\item C、D，$n=1023$：\Mtwo{}($w=1/16$) 慢于 BoomerAMG——\textbf{成立}（3.1--3.4 倍）；最快的环宽在 1.5 倍以内——\textbf{不成立}（1.44--1.82 倍）。
\item C、D：\Mtwo{} 停止时代数误差大于 BoomerAMG——\textbf{成立}（64/64）；$k_8\le14$——\textbf{成立}（最大 11）。
\end{enumerate}

\section{三维结果}
\IfFileExists{tables_3d.tex}{\begin{table}[H]\centering\small
\caption{三维：PCG 迭代次数（$\chi\in\{10^2,10^4,10^6,10^8\}$ 的最小--最大值；$^{\times k}$ 表示有 $k$ 个 $\chi$ 在 500 次内未收敛）}\label{tab:its3d}
\begin{adjustbox}{max width=\textwidth}
\begin{tabular}{llcccccc}
\toprule
算例 & $n$ & 传统 LGF & GMG & BoomerAMG & 方法二 & M2 $w{=}\frac1{32}$ & M2 $w{=}\frac1{64}$ \\
\midrule
A: $\chi^x$ & 31 & 68--500$^{\times2}$ & 8--9 & 5--7 & 2--3 &  &  \\
 & 63 & 97--500$^{\times3}$ & 9 & 7 & 2--3 &  &  \\
 & 127 & 111--500$^{\times3}$ & 9 & 7--9 & 2--3 &  &  \\
\midrule
B: 阶跃层 & 31 & 20--48 & 9 & 7--8 & 1--3 &  &  \\
 & 63 & 19--59 & 9 & 9 & 1--3 &  &  \\
 & 127 & 19--62 & 9 & 10--11 & 1--3 &  &  \\
\midrule
C: $\chi^x$+规则夹杂 & 31 & 500$^{\times4}$ & 57--62 & 11 &  & 12 & 12 \\
 & 63 & 500$^{\times4}$ & 75--89 & 12--14 &  & 11 & 15--16 \\
 & 127 & 500$^{\times4}$ & 89--103 & 14--15 &  & 10 & 13--14 \\
\midrule
D: $\chi^x$+随机夹杂 & 31 & 500$^{\times4}$ & 70--88 & 13--14 &  & 12--13 & 12--13 \\
 & 63 & 500$^{\times4}$ & 100--107 & 15--16 &  & 12 & 17 \\
 & 127 & 500$^{\times4}$ & 123--135 & 15--17 &  & 11 & 16 \\
\bottomrule
\end{tabular}
\end{adjustbox}
\end{table}

\begin{table}[H]\centering\small
\caption{三维：误差 $\|u_k-u_h\|/\|u_h\|\le10^{-8}$ 所需迭代数（各 $\chi$ 的范围；--：300 次内未达到）}\label{tab:k83d}
\begin{adjustbox}{max width=\textwidth}
\begin{tabular}{llcccccc}
\toprule
算例 & $n$ & 传统 LGF & GMG & BoomerAMG & 方法二 & M2 $w{=}\frac1{32}$ & M2 $w{=}\frac1{64}$ \\
\midrule
A: $\chi^x$ & 31 & 54--244 & 6--8 & 5 & 2--3 &  &  \\
 & 63 & 65 & 6--7 & 5--6 & 2--3 &  &  \\
 & 127 & 70 & 6--7 & 6--8 & 2--3 &  &  \\
\midrule
B: 阶跃层 & 31 & 12--80 & 6 & 6 & 3 &  &  \\
 & 63 & 12--85 & 6 & 7--8 & 3 &  &  \\
 & 127 & 12--87 & 5 & 8--9 & 3 &  &  \\
\midrule
C: $\chi^x$+规则夹杂 & 31 & -- & 64--67 & 10--11 &  & 13--14 & 13--14 \\
 & 63 & -- & 84--86 & 11--14 &  & 12--14 & 15--18 \\
 & 127 & -- & 99--100 & 13--14 &  & 10--13 & 14--16 \\
\midrule
D: $\chi^x$+随机夹杂 & 31 & -- & 80--86 & 12--14 &  & 13--15 & 13--15 \\
 & 63 & -- & 104--110 & 13--16 &  & 12--13 & 17--19 \\
 & 127 & -- & 123--129 & 15 &  & 11--12 & 16--18 \\
\bottomrule
\end{tabular}
\end{adjustbox}
\end{table}

\begin{table}[H]\centering\small
\caption{三维：按残差规则 $\|r\|/\|f\|\le10^{-10}$ 停止时的代数误差 $\|u-u_h\|/\|u_h\|$（各 $\chi$ 中的最大值）}\label{tab:err3d}
\begin{adjustbox}{max width=\textwidth}
\begin{tabular}{llcccccc}
\toprule
算例 & $n$ & 传统 LGF & GMG & BoomerAMG & 方法二 & M2 $w{=}\frac1{32}$ & M2 $w{=}\frac1{64}$ \\
\midrule
A: $\chi^x$ & 31 & $2.4{\times}10^{-2}$ & $6.1{\times}10^{-9}$ & $1.7{\times}10^{-10}$ & $3.3{\times}10^{-9}$ &  &  \\
 & 63 & $5.8{\times}10^{-2}$ & $7.1{\times}10^{-11}$ & $1.5{\times}10^{-10}$ & $3.3{\times}10^{-9}$ &  &  \\
 & 127 & $8.4{\times}10^{-2}$ & $3.2{\times}10^{-11}$ & $1.5{\times}10^{-10}$ & $3.3{\times}10^{-9}$ &  &  \\
\midrule
B: 阶跃层 & 31 & $7.3{\times}10^{-7}$ & $3.3{\times}10^{-12}$ & $4.1{\times}10^{-10}$ & $2.0{\times}10^{-4}$ &  &  \\
 & 63 & $7.5{\times}10^{-7}$ & $8.9{\times}10^{-13}$ & $1.6{\times}10^{-10}$ & $2.0{\times}10^{-4}$ &  &  \\
 & 127 & $7.5{\times}10^{-7}$ & $3.4{\times}10^{-13}$ & $3.9{\times}10^{-10}$ & $2.0{\times}10^{-4}$ &  &  \\
\midrule
C: $\chi^x$+规则夹杂 & 31 & $3.3{\times}10^{-1}$ & $1.0{\times}10^{-7}$ & $3.3{\times}10^{-9}$ &  & $2.7{\times}10^{-7}$ & $2.7{\times}10^{-7}$ \\
 & 63 & $4.1{\times}10^{-1}$ & $2.2{\times}10^{-7}$ & $1.5{\times}10^{-8}$ &  & $4.5{\times}10^{-7}$ & $3.1{\times}10^{-7}$ \\
 & 127 & $4.4{\times}10^{-1}$ & $2.9{\times}10^{-7}$ & $2.4{\times}10^{-9}$ &  & $2.5{\times}10^{-7}$ & $1.2{\times}10^{-7}$ \\
\midrule
D: $\chi^x$+随机夹杂 & 31 & $5.4{\times}10^{-1}$ & $3.0{\times}10^{-7}$ & $6.6{\times}10^{-9}$ &  & $7.0{\times}10^{-8}$ & $7.0{\times}10^{-8}$ \\
 & 63 & $6.3{\times}10^{-1}$ & $6.0{\times}10^{-8}$ & $1.3{\times}10^{-8}$ &  & $1.2{\times}10^{-8}$ & $7.9{\times}10^{-8}$ \\
 & 127 & $6.6{\times}10^{-1}$ & $9.9{\times}10^{-9}$ & $4.4{\times}10^{-9}$ &  & $3.2{\times}10^{-8}$ & $1.7{\times}10^{-7}$ \\
\bottomrule
\end{tabular}
\end{adjustbox}
\end{table}

\begin{table}[H]\centering\small
\caption{三维：总时间 setup+solve（$\chi=10^8$，单线程，ms）}\label{tab:time3d}
\begin{adjustbox}{max width=\textwidth}
\begin{tabular}{llcccccc}
\toprule
算例 & $n$ & 传统 LGF & GMG & BoomerAMG & 方法二 & M2 $w{=}\frac1{32}$ & M2 $w{=}\frac1{64}$ \\
\midrule
A: $\chi^x$ & 31 & 745 & 33 & 100 & 5 &  &  \\
 & 63 & 6.52\,s & 257 & 1.23\,s & 41 &  &  \\
 & 127 & 47.96\,s & 2.62\,s & 8.86\,s & 335 &  &  \\
\midrule
B: 阶跃层 & 31 & 69 & 31 & 119 & 4 &  &  \\
 & 63 & 698 & 265 & 1.21\,s & 29 &  &  \\
 & 127 & 5.14\,s & 2.28\,s & 12.80\,s & 229 &  &  \\
\midrule
C: $\chi^x$+规则夹杂 & 31 & 688 & 137 & 134 &  & 48 & 67 \\
 & 63 & 6.25\,s & 1.49\,s & 1.37\,s &  & 503 & 429 \\
 & 127 & 44.08\,s & 13.50\,s & 10.09\,s &  & 21.21\,s & 9.25\,s \\
\midrule
D: $\chi^x$+随机夹杂 & 31 & 710 & 164 & 155 &  & 41 & 42 \\
 & 63 & 6.29\,s & 1.93\,s & 1.66\,s &  & 369 & 376 \\
 & 127 & 43.00\,s & 19.73\,s & 11.66\,s &  & 9.80\,s & 6.01\,s \\
\bottomrule
\end{tabular}
\end{adjustbox}
\end{table}

\begin{table}[H]\centering\small
\caption{三维：误差降到 $10^{-8}$ 的时间 setup$+k_8\times$单步时间（$\chi=10^8$，ms）}\label{tab:t83d}
\begin{adjustbox}{max width=\textwidth}
\begin{tabular}{llcccccc}
\toprule
算例 & $n$ & 传统 LGF & GMG & BoomerAMG & 方法二 & M2 $w{=}\frac1{32}$ & M2 $w{=}\frac1{64}$ \\
\midrule
A: $\chi^x$ & 31 & -- & 33 & 100 & 5 &  &  \\
 & 63 & -- & 219 & 1.15\,s & 41 &  &  \\
 & 127 & -- & 2.25\,s & 8.40\,s & 335 &  &  \\
\midrule
B: 阶跃层 & 31 & 115 & 24 & 112 & 10 &  &  \\
 & 63 & 1.12\,s & 205 & 1.15\,s & 78 &  &  \\
 & 127 & 7.98\,s & 1.61\,s & 11.56\,s & 611 &  &  \\
\midrule
C: $\chi^x$+规则夹杂 & 31 & -- & 153 & 127 &  & 54 & 75 \\
 & 63 & -- & 1.66\,s & 1.43\,s &  & 576 & 489 \\
 & 127 & -- & 14.95\,s & 9.66\,s &  & 22.29\,s & 9.99\,s \\
\midrule
D: $\chi^x$+随机夹杂 & 31 & -- & 188 & 155 &  & 49 & 50 \\
 & 63 & -- & 2.07\,s & 1.73\,s &  & 390 & 412 \\
 & 127 & -- & 19.73\,s & 11.66\,s &  & 10.07\,s & 6.39\,s \\
\bottomrule
\end{tabular}
\end{adjustbox}
\end{table}

\begin{table}[H]\centering\small
\caption{三维，算例 A（$n=127$）：相对连续解的总误差 $\|u-u^*\|/\|u^*\|$}\label{tab:aerr三维}
\begin{tabular}{lcccc}
\toprule
$\chi$ & 传统 LGF & GMG & BoomerAMG & 方法二 \\
\midrule
$10^{2}$ & $6.3{\times}10^{-5}$ & $6.3{\times}10^{-5}$ & $6.3{\times}10^{-5}$ & $6.3{\times}10^{-5}$ \\
$10^{4}$ & $4.8{\times}10^{-4}$ & $4.8{\times}10^{-4}$ & $4.8{\times}10^{-4}$ & $4.8{\times}10^{-4}$ \\
$10^{6}$ & $9.3{\times}10^{-3}$ & $1.3{\times}10^{-3}$ & $1.3{\times}10^{-3}$ & $1.3{\times}10^{-3}$ \\
$10^{8}$ & $8.5{\times}10^{-2}$ & $2.5{\times}10^{-3}$ & $2.5{\times}10^{-3}$ & $2.5{\times}10^{-3}$ \\
\bottomrule
\end{tabular}
\end{table}}{}

\begin{figure}[H]\centering
\includegraphics[width=\textwidth]{fig3d_iters_vs_n.png}
\caption{三维：$\chi=10^8$ 时迭代次数随 $n$ 的变化。}\label{fig:3dits}
\end{figure}
\begin{figure}[H]\centering
\includegraphics[width=\textwidth]{fig3d_error_traces.png}
\caption{三维：$n=127$、$\chi=10^8$ 时误差随迭代的下降。}\label{fig:3dtrace}
\end{figure}
\begin{figure}[H]\centering
\includegraphics[width=\textwidth]{fig3d_time_to_1e-8_vs_n.png}
\caption{三维：达到误差 $10^{-8}$ 的时间（$\chi=10^8$）。}\label{fig:3dtime}
\end{figure}

\subsection{三维分析}
\paragraph{迭代次数。}A、B 上\Mtwo{} 仍是 1--3 次。BoomerAMG 在三维随 $n$ 略增（A 5--9 次，B 7--11 次），
GMG（$\omega=6/7$）稳定在 8--9 次；传统 LGF 在 A 上 $\chi=10^2$ 要 68--111 次，$\chi\ge10^4$（$n\ge63$）时 500 次内不收敛，
在 B 上 19--62 次。C、D 上\Mtwo{}（$w=1/32$）10--13 次，$n=127$ 时只要 10--11 次；$w=1/64$ 要 12--17 次；
BoomerAMG 从 $n=31$ 的 11--14 次增加到 $n=127$ 的 14--17 次；GMG 57--135 次且随 $n$ 增长；传统 LGF 全部不收敛。

\paragraph{精度。}A 上所有收敛方法总误差前 3 位相同（$n=127$、$\chi=10^8$：$2.53\times10^{-3}$；三维网格较粗，离散误差比二维大）；
传统 LGF 在 $\chi\ge10^6$ 时没收敛，误差最多大 34 倍。按残差停止时，\Mtwo{} 在 C、D 上的代数误差为 $1.4\times10^{-9}$--$4.5\times10^{-7}$，
BoomerAMG 为 $1.8\times10^{-10}$--$1.5\times10^{-8}$（48 组中 46 组\Mtwo{} 更大）；B 上 $\chi=10^8$ 时\Mtwo{} 只迭代 1 次，误差 $2\times10^{-4}$，
与二维相同的停止规则问题。按误差计数，把误差降到 $10^{-8}$：\Mtwo{}（$w=1/32$）10--15 次，BoomerAMG 10--16 次，GMG 64--129 次。
图~\ref{fig:3dtrace} 中，$w=1/32$ 的误差曲线下降最快，BoomerAMG 次之。

\paragraph{时间（$n=127$，约 $2\times10^6$ 个未知量）。}A、B：\Mtwo{} 0.23--0.47 s，BoomerAMG 7.1--13.5 s（其中 setup 3.4--6.2 s），
GMG 2.2--2.9 s，传统 LGF 1.9--48 s；\Mtwo{} 是 BoomerAMG 的 $\le0.055$ 倍、GMG 的 0.10--0.21 倍。
三维时 GMG 比 BoomerAMG 快 3--6 倍（BoomerAMG 在三维的 setup 和算子复杂度更高）。
C、D：三维的夹杂是彼此分开的小球，膨胀后 LU 区域只占网格的 6--15\%（二维 $w=1/16$ 时为 48--57\%），分成 23--29 个小区域分别分解。
$w=1/64$ 时 C 上 9.2--9.9 s，与 BoomerAMG（9.3--10.1 s）持平；D 上 5.5--6.0 s，BoomerAMG 11.0--11.7 s，快约 2 倍。
$w=1/32$ 迭代更少，但三维 LU 的填充增长很快（C 上 setup 从 $n=63$ 的 0.24 s 增加到 $n=127$ 的 17 s），总时间 20--21 s（C）、9.2--9.8 s（D）。
GMG 13.5--21 s，传统 LGF 42--47 s 且不收敛。

\paragraph{预测检验（三维，锁定方案 \texttt{06cc1590}）。}
\begin{enumerate}
\item A、B：\Mtwo{} $\le3$ 次；BoomerAMG、GMG $\le15$ 次；LGF 在 A 上 $\chi\ge10^6$ 不收敛——\textbf{成立}。
\item C、D：\Mtwo{}($w=1/32$) $\le15$ 次且从 $n=31$ 到 127 变化 $\le3$；$n=127$ 时 BoomerAMG 迭代多于\Mtwo{}；GMG $\ge50$——\textbf{成立}。
\item $n=127$ 时间：A、B 上\Mtwo{} $\le$ BoomerAMG 的 1/3——\textbf{成立}（$\le0.055$）；C、D 上最快环宽在 BoomerAMG 的 1.3 倍以内——\textbf{成立}；
  “不分胜负”——C 上成立（0.92--1.03），D 上\textbf{不成立}（\Mtwo{} 快约 2 倍）；$w=1/32$ 慢于 $w=1/64$——\textbf{成立}。
\item A：收敛方法总误差 3 位相同——\textbf{成立}。
\item C、D：\Mtwo{} 停止时代数误差大于 BoomerAMG——\textbf{基本成立}（$w=1/32$ 46/48，$w=1/64$ 48/48）；$k_8\le18$——\textbf{成立}（最大 15）。
\end{enumerate}
锁定前的冒烟测试已看到 $n=63$ 与 $n=127$（仅 C、$\chi=10^6$）的部分结果，见方案文件；因此第 3 条的时间预测不是完全盲测。

\section{与其他方法的区别（文献对比）}
详细内容见 \texttt{FD\_LGF/literature.md}。下面按类别概括；多数文献只读了摘要或二手引用。
\begin{itemize}
\item \textbf{传统 LGF / FFT 参考介质}（Xia；Moulinec--Suquet 1998；Zeman 等 2010）：用一个常系数算子的快速逆，
  迭代数随对比度增长（CG 加速后约 $\sqrt\chi$）。\Mtwo{} 把对比度放进两个精确部件：分层部分在 $G$ 里、夹杂在 $Q$ 里。
\item \textbf{可分离 / 快速对角化预条件}（Concus--Golub 1973；Lynch--Rice--Thomas 1964；Huynh--Ito--Toivanen）：
  只对 $a_1(x)a_2(y)$ 精确；\Mtwo{} 在此基础上加入夹杂附近的精确局部修正。
\item \textbf{电容矩阵法}（Buzbee 等 1971；Proskurowski--Widlund 1976）与 \textbf{fat boundary method}（Maury 2001）：
  全区域快速解法 $+$ 孔洞或边界附近的小问题。后者结构上最接近\Mtwo{}，区别是：它处理孔洞（局部问题加 Dirichlet 条件），
  用带松弛参数 $\theta$ 的不动点迭代耦合；\Mtwo{} 处理有限系数、可高可低的夹杂，纯代数（直接取 $A$ 的子块），
  是对称正定预条件，放在 CG 中，无需调参，且全局快速解法针对变系数分层背景。
\item \textbf{两层区域分解 / BNN / deflation / GenEO}（Mandel 1993；Tang 等 2009；Graham--Lechner--Scheichl 2007；Spillane 等 2014）：
  代数形式相同，但其中 $Q$ 是低维粗空间，必须随系数构造才能对高对比度稳健；\Mtwo{} 中 $Q$ 是局部精确解，
  $G$ 是全局、全尺度的精确背景逆，不需要粗空间。用 Xu (1992) 的子空间校正语言，\Mtwo{} 是两个子空间的对称逐次校正。
\item \textbf{AGKS}（Aksoylu--Graham--Klie--Scheichl 2008）：按高/低系数分块；\Mtwo{} 按“背景/特征”分块，背景自身可以有 $10^8$ 的对比度。
\item \textbf{GMG / BoomerAMG}（Alcouffe 等 1981；Henson--Yang 2002）：GMG 的几何插值在系数跳跃处退化（本文 C、D 上二维 48--114 次、三维 57--135 次）；
  BoomerAMG 是黑箱，自动适应系数并可并行。\Mtwo{} 是“灰箱”：需要笛卡尔网格、可分离背景和可识别的夹杂。
\end{itemize}

\section{结论与局限}
\begin{enumerate}
\item \textbf{教授的 $a=\chi^x$ 和分层系数：\Mtwo{} 最好。}二维、三维都只需 1--3 次迭代，与 $\chi\le10^8$、$n$ 都无关；
  setup 只是一次三对角分解；时间比 BoomerAMG 少一个数量级以上（三维），比教科书参数的 GMG 少 2.4--10 倍。
  这时\Mtwo{} 就是方法一（精确的分层 LGF），因为没有夹杂。
\item \textbf{分层背景加夹杂：迭代数稳健，时间看 LU 规模。}\Mtwo{} 的迭代数与 $\chi$、$n$ 无关，而 BoomerAMG 在三维随 $n$ 增长、
  GMG 随 $n$ 明显增长。二维时夹杂区（含环）占网格一半，LU 太贵，\Mtwo{} 比 BoomerAMG 慢 1.4--3.4 倍；
  三维时夹杂区只占 6--15\%，\Mtwo{} 与 BoomerAMG 持平或快约 2 倍。窄环（$w=1/64$）通常更省时间。
\item \textbf{精度：}收敛时所有方法的离散误差相同，差别只在代数误差。残差停止规则会低估\Mtwo{} 在低系数区的误差，
  高对比度下应按误差（或预条件残差）停止，或把容差收紧到 $10^{-13}$。
\item \textbf{传统 LGF}（常系数参考）在 $\chi\ge10^4$ 时基本不可用，这正是\Mtwo{} 要解决的问题。
\end{enumerate}
\paragraph{局限。}(i) \Mtwo{} 需要结构信息：笛卡尔网格、可分离的背景、可识别的夹杂；背景不可分离时需要先拟合。
(ii) 局部解用稀疏直接 LU，夹杂体积大或网格更细时 setup 超线性增长（三维 $w=1/32$ 已明显）；不精确的局部解会破坏收敛（H9）。
(iii) 本文全部是单线程 Python 实现，BoomerAMG 有成熟的 MPI 并行版本，\Mtwo{} 还没有。
(iv) 计时时机器上有其他程序（load 2--12），迭代次数和误差是主要指标，时间只作相对比较。

\appendix
\section{完整结果表}
\begingroup\footnotesize\setlength{\tabcolsep}{3pt}
\begin{longtable}{lrrlrcccrrr}
\caption{二维 全部结果。err：按残差规则停止时相对 $u_h$ 的代数误差；$k_8$：误差降到 $10^{-8}$ 所需迭代数；时间单位 ms（单线程）。}\\
\toprule 算例 & $n$ & $\chi$ & 方法 & 迭代 & err & 第 5 步误差 & $k_8$ & setup & solve & 合计 \\ \midrule
\endfirsthead \toprule 算例 & $n$ & $\chi$ & 方法 & 迭代 & err & 第 5 步误差 & $k_8$ & setup & solve & 合计 \\ \midrule \endhead
A & 127 & $10^{2}$ & 传统 LGF & 110 & $5.5{\times}10^{-13}$ & $1.4{\times}10^{-1}$ & 68 & 0 & 59 & 59 \\
A & 127 & $10^{2}$ & BoomerAMG & 7 & $2.0{\times}10^{-13}$ & $4.1{\times}10^{-10}$ & 5 & 9 & 10 & 19 \\
A & 127 & $10^{2}$ & M2 $w{=}\frac1{16}$ & 3 & $3.7{\times}10^{-12}$ & $1.9{\times}10^{-14}$ & 3 & 2 & 3 & 5 \\
A & 127 & $10^{2}$ & GMG & 7 & $4.0{\times}10^{-13}$ & $4.7{\times}10^{-9}$ & 5 & 6 & 8 & 14 \\
A & 127 & $10^{4}$ & 传统 LGF & 781 & $1.5{\times}10^{-11}$ & $4.7{\times}10^{-1}$ & -- & 0 & 415 & 415 \\
A & 127 & $10^{4}$ & BoomerAMG & 7 & $8.2{\times}10^{-13}$ & $2.2{\times}10^{-9}$ & 5 & 9 & 9 & 17 \\
A & 127 & $10^{4}$ & M2 $w{=}\frac1{16}$ & 3 & $3.6{\times}10^{-13}$ & $2.6{\times}10^{-14}$ & 2 & 2 & 3 & 5 \\
A & 127 & $10^{4}$ & GMG & 7 & $3.3{\times}10^{-12}$ & $4.0{\times}10^{-8}$ & 6 & 6 & 8 & 14 \\
A & 127 & $10^{6}$ & 传统 LGF & 1000$^\times$ & $8.0{\times}10^{-4}$ & $6.4{\times}10^{-1}$ & -- & 0 & 542 & 542 \\
A & 127 & $10^{6}$ & BoomerAMG & 7 & $2.1{\times}10^{-12}$ & $4.4{\times}10^{-9}$ & 5 & 8 & 9 & 17 \\
A & 127 & $10^{6}$ & M2 $w{=}\frac1{16}$ & 2 & $1.4{\times}10^{-9}$ & $7.9{\times}10^{-14}$ & 2 & 1 & 2 & 4 \\
A & 127 & $10^{6}$ & GMG & 7 & $4.7{\times}10^{-11}$ & $1.1{\times}10^{-7}$ & 6 & 6 & 8 & 14 \\
A & 127 & $10^{8}$ & 传统 LGF & 1000$^\times$ & $2.8{\times}10^{-2}$ & $7.4{\times}10^{-1}$ & -- & 1 & 542 & 542 \\
A & 127 & $10^{8}$ & BoomerAMG & 7 & $3.5{\times}10^{-12}$ & $6.1{\times}10^{-9}$ & 5 & 8 & 10 & 18 \\
A & 127 & $10^{8}$ & M2 $w{=}\frac1{16}$ & 2 & $4.2{\times}10^{-10}$ & $2.3{\times}10^{-14}$ & 2 & 1 & 2 & 4 \\
A & 127 & $10^{8}$ & GMG & 7 & $3.2{\times}10^{-10}$ & $6.5{\times}10^{-7}$ & 7 & 6 & 9 & 16 \\
A & 255 & $10^{2}$ & 传统 LGF & 117 & $3.6{\times}10^{-13}$ & $1.4{\times}10^{-1}$ & 71 & 1 & 258 & 260 \\
A & 255 & $10^{2}$ & BoomerAMG & 7 & $1.8{\times}10^{-13}$ & $5.1{\times}10^{-10}$ & 5 & 35 & 37 & 71 \\
A & 255 & $10^{2}$ & M2 $w{=}\frac1{16}$ & 3 & $3.8{\times}10^{-12}$ & $5.9{\times}10^{-14}$ & 3 & 6 & 11 & 17 \\
A & 255 & $10^{2}$ & GMG & 7 & $1.6{\times}10^{-13}$ & $2.7{\times}10^{-9}$ & 5 & 20 & 31 & 51 \\
A & 255 & $10^{4}$ & 传统 LGF & 900 & $8.2{\times}10^{-12}$ & $4.7{\times}10^{-1}$ & -- & 1 & 2015 & 2016 \\
A & 255 & $10^{4}$ & BoomerAMG & 7 & $8.3{\times}10^{-13}$ & $1.5{\times}10^{-9}$ & 5 & 36 & 37 & 73 \\
A & 255 & $10^{4}$ & M2 $w{=}\frac1{16}$ & 3 & $3.7{\times}10^{-13}$ & $6.7{\times}10^{-14}$ & 2 & 6 & 11 & 17 \\
A & 255 & $10^{4}$ & GMG & 8 & $1.7{\times}10^{-13}$ & $3.3{\times}10^{-8}$ & 6 & 19 & 33 & 52 \\
A & 255 & $10^{6}$ & 传统 LGF & 1000$^\times$ & $1.4{\times}10^{-3}$ & $6.5{\times}10^{-1}$ & -- & 2 & 2350 & 2352 \\
A & 255 & $10^{6}$ & BoomerAMG & 7 & $5.1{\times}10^{-12}$ & $8.9{\times}10^{-9}$ & 5 & 36 & 41 & 76 \\
A & 255 & $10^{6}$ & M2 $w{=}\frac1{16}$ & 2 & $1.4{\times}10^{-9}$ & $7.7{\times}10^{-14}$ & 2 & 5 & 8 & 13 \\
A & 255 & $10^{6}$ & GMG & 8 & $6.5{\times}10^{-13}$ & $8.8{\times}10^{-8}$ & 6 & 20 & 36 & 56 \\
A & 255 & $10^{8}$ & 传统 LGF & 1000$^\times$ & $3.7{\times}10^{-2}$ & $7.4{\times}10^{-1}$ & -- & 2 & 2359 & 2361 \\
A & 255 & $10^{8}$ & BoomerAMG & 7 & $4.1{\times}10^{-12}$ & $6.4{\times}10^{-9}$ & 5 & 37 & 40 & 77 \\
A & 255 & $10^{8}$ & M2 $w{=}\frac1{16}$ & 2 & $4.2{\times}10^{-10}$ & $3.0{\times}10^{-14}$ & 2 & 6 & 9 & 15 \\
A & 255 & $10^{8}$ & GMG & 8 & $3.0{\times}10^{-12}$ & $5.1{\times}10^{-7}$ & 7 & 20 & 35 & 55 \\
A & 511 & $10^{2}$ & 传统 LGF & 123 & $1.8{\times}10^{-13}$ & $1.4{\times}10^{-1}$ & 73 & 7 & 1350 & 1356 \\
A & 511 & $10^{2}$ & BoomerAMG & 7 & $1.8{\times}10^{-13}$ & $6.2{\times}10^{-10}$ & 5 & 152 & 164 & 316 \\
A & 511 & $10^{2}$ & M2 $w{=}\frac1{16}$ & 3 & $3.8{\times}10^{-12}$ & $5.6{\times}10^{-14}$ & 3 & 25 & 54 & 79 \\
A & 511 & $10^{2}$ & GMG & 7 & $8.7{\times}10^{-14}$ & $2.0{\times}10^{-9}$ & 5 & 79 & 126 & 205 \\
A & 511 & $10^{4}$ & 传统 LGF & 979 & $4.1{\times}10^{-12}$ & $4.7{\times}10^{-1}$ & -- & 8 & 10638 & 10646 \\
A & 511 & $10^{4}$ & BoomerAMG & 7 & $1.2{\times}10^{-12}$ & $1.9{\times}10^{-9}$ & 5 & 169 & 166 & 335 \\
A & 511 & $10^{4}$ & M2 $w{=}\frac1{16}$ & 3 & $5.1{\times}10^{-13}$ & $2.9{\times}10^{-13}$ & 2 & 24 & 57 & 82 \\
A & 511 & $10^{4}$ & GMG & 8 & $3.1{\times}10^{-13}$ & $3.1{\times}10^{-8}$ & 6 & 78 & 143 & 221 \\
A & 511 & $10^{6}$ & 传统 LGF & 1000$^\times$ & $1.9{\times}10^{-3}$ & $6.5{\times}10^{-1}$ & -- & 7 & 10717 & 10725 \\
A & 511 & $10^{6}$ & BoomerAMG & 7 & $4.1{\times}10^{-12}$ & $7.8{\times}10^{-9}$ & 5 & 151 & 167 & 317 \\
A & 511 & $10^{6}$ & M2 $w{=}\frac1{16}$ & 2 & $1.4{\times}10^{-9}$ & $1.4{\times}10^{-13}$ & 2 & 26 & 45 & 72 \\
A & 511 & $10^{6}$ & GMG & 8 & $5.5{\times}10^{-13}$ & $8.4{\times}10^{-8}$ & 6 & 78 & 131 & 208 \\
A & 511 & $10^{8}$ & 传统 LGF & 1000$^\times$ & $4.2{\times}10^{-2}$ & $7.4{\times}10^{-1}$ & -- & 7 & 11062 & 11069 \\
A & 511 & $10^{8}$ & BoomerAMG & 7 & $5.9{\times}10^{-12}$ & $6.9{\times}10^{-9}$ & 5 & 146 & 160 & 306 \\
A & 511 & $10^{8}$ & M2 $w{=}\frac1{16}$ & 2 & $4.2{\times}10^{-10}$ & $5.0{\times}10^{-14}$ & 2 & 25 & 37 & 62 \\
A & 511 & $10^{8}$ & GMG & 8 & $1.9{\times}10^{-12}$ & $4.4{\times}10^{-7}$ & 6 & 77 & 142 & 219 \\
A & 1023 & $10^{2}$ & 传统 LGF & 127 & $3.2{\times}10^{-13}$ & $1.4{\times}10^{-1}$ & 73 & 25 & 5717 & 5742 \\
A & 1023 & $10^{2}$ & BoomerAMG & 8 & $3.1{\times}10^{-13}$ & $8.1{\times}10^{-10}$ & 5 & 588 & 720 & 1308 \\
A & 1023 & $10^{2}$ & M2 $w{=}\frac1{16}$ & 3 & $3.7{\times}10^{-12}$ & $3.0{\times}10^{-13}$ & 3 & 93 & 247 & 340 \\
A & 1023 & $10^{2}$ & GMG & 7 & $3.4{\times}10^{-13}$ & $1.6{\times}10^{-9}$ & 5 & 325 & 605 & 930 \\
A & 1023 & $10^{4}$ & 传统 LGF & 1000$^\times$ & $4.6{\times}10^{-12}$ & $4.7{\times}10^{-1}$ & -- & 23 & 49079 & 49102 \\
A & 1023 & $10^{4}$ & BoomerAMG & 7 & $1.1{\times}10^{-12}$ & $1.7{\times}10^{-9}$ & 5 & 574 & 593 & 1166 \\
A & 1023 & $10^{4}$ & M2 $w{=}\frac1{16}$ & 3 & $1.3{\times}10^{-12}$ & $1.1{\times}10^{-12}$ & 2 & 89 & 179 & 268 \\
A & 1023 & $10^{4}$ & GMG & 8 & $1.1{\times}10^{-12}$ & $2.7{\times}10^{-8}$ & 6 & 320 & 644 & 964 \\
A & 1023 & $10^{6}$ & 传统 LGF & 1000$^\times$ & $2.1{\times}10^{-3}$ & $6.5{\times}10^{-1}$ & -- & 25 & 48157 & 48183 \\
A & 1023 & $10^{6}$ & BoomerAMG & 7 & $5.6{\times}10^{-12}$ & $8.3{\times}10^{-9}$ & 5 & 561 & 577 & 1138 \\
A & 1023 & $10^{6}$ & M2 $w{=}\frac1{16}$ & 2 & $1.4{\times}10^{-9}$ & $4.2{\times}10^{-13}$ & 2 & 112 & 215 & 327 \\
A & 1023 & $10^{6}$ & GMG & 8 & $4.6{\times}10^{-13}$ & $8.2{\times}10^{-8}$ & 6 & 312 & 604 & 916 \\
A & 1023 & $10^{8}$ & 传统 LGF & 1000$^\times$ & $4.5{\times}10^{-2}$ & $7.4{\times}10^{-1}$ & -- & 33 & 46268 & 46301 \\
A & 1023 & $10^{8}$ & BoomerAMG & 7 & $1.1{\times}10^{-11}$ & $1.0{\times}10^{-8}$ & 5 & 567 & 569 & 1136 \\
A & 1023 & $10^{8}$ & M2 $w{=}\frac1{16}$ & 2 & $4.2{\times}10^{-10}$ & $2.9{\times}10^{-13}$ & 2 & 89 & 162 & 251 \\
A & 1023 & $10^{8}$ & GMG & 8 & $1.3{\times}10^{-12}$ & $3.9{\times}10^{-7}$ & 6 & 312 & 587 & 899 \\
B & 127 & $10^{2}$ & 传统 LGF & 15 & $2.3{\times}10^{-12}$ & $3.6{\times}10^{-4}$ & 11 & 0 & 8 & 9 \\
B & 127 & $10^{2}$ & BoomerAMG & 8 & $3.3{\times}10^{-13}$ & $2.2{\times}10^{-8}$ & 6 & 10 & 13 & 23 \\
B & 127 & $10^{2}$ & M2 $w{=}\frac1{16}$ & 3 & $3.6{\times}10^{-11}$ & $1.2{\times}10^{-14}$ & 3 & 1 & 3 & 4 \\
B & 127 & $10^{2}$ & GMG & 7 & $7.7{\times}10^{-13}$ & $5.5{\times}10^{-10}$ & 5 & 6 & 9 & 14 \\
B & 127 & $10^{4}$ & 传统 LGF & 29 & $8.2{\times}10^{-10}$ & $3.8{\times}10^{-4}$ & 18 & 0 & 17 & 17 \\
B & 127 & $10^{4}$ & BoomerAMG & 8 & $4.4{\times}10^{-12}$ & $4.1{\times}10^{-8}$ & 6 & 9 & 11 & 20 \\
B & 127 & $10^{4}$ & M2 $w{=}\frac1{16}$ & 3 & $3.1{\times}10^{-11}$ & $1.1{\times}10^{-14}$ & 3 & 1 & 3 & 4 \\
B & 127 & $10^{4}$ & GMG & 7 & $3.1{\times}10^{-12}$ & $1.7{\times}10^{-9}$ & 5 & 6 & 9 & 15 \\
B & 127 & $10^{6}$ & 传统 LGF & 51 & $8.2{\times}10^{-10}$ & $1.3{\times}10^{-3}$ & 43 & 0 & 28 & 28 \\
B & 127 & $10^{6}$ & BoomerAMG & 8 & $6.5{\times}10^{-12}$ & $4.1{\times}10^{-8}$ & 6 & 10 & 12 & 21 \\
B & 127 & $10^{6}$ & M2 $w{=}\frac1{16}$ & 2 & $2.2{\times}10^{-7}$ & $1.1{\times}10^{-14}$ & 3 & 1 & 2 & 3 \\
B & 127 & $10^{6}$ & GMG & 7 & $3.3{\times}10^{-12}$ & $1.7{\times}10^{-9}$ & 5 & 6 & 9 & 15 \\
B & 127 & $10^{8}$ & 传统 LGF & 45 & $6.0{\times}10^{-7}$ & $9.0{\times}10^{-2}$ & 69 & 0 & 26 & 26 \\
B & 127 & $10^{8}$ & BoomerAMG & 8 & $6.5{\times}10^{-12}$ & $4.1{\times}10^{-8}$ & 6 & 9 & 12 & 21 \\
B & 127 & $10^{8}$ & M2 $w{=}\frac1{16}$ & 1 & $3.7{\times}10^{-4}$ & $1.1{\times}10^{-14}$ & 3 & 1 & 2 & 3 \\
B & 127 & $10^{8}$ & GMG & 7 & $3.3{\times}10^{-12}$ & $1.7{\times}10^{-9}$ & 5 & 6 & 9 & 15 \\
B & 255 & $10^{2}$ & 传统 LGF & 15 & $2.3{\times}10^{-12}$ & $3.6{\times}10^{-4}$ & 11 & 1 & 36 & 37 \\
B & 255 & $10^{2}$ & BoomerAMG & 8 & $3.9{\times}10^{-13}$ & $1.5{\times}10^{-8}$ & 6 & 36 & 44 & 80 \\
B & 255 & $10^{2}$ & M2 $w{=}\frac1{16}$ & 3 & $3.6{\times}10^{-11}$ & $1.2{\times}10^{-14}$ & 3 & 6 & 11 & 16 \\
B & 255 & $10^{2}$ & GMG & 7 & $5.4{\times}10^{-13}$ & $5.0{\times}10^{-11}$ & 5 & 19 & 32 & 51 \\
B & 255 & $10^{4}$ & 传统 LGF & 31 & $8.2{\times}10^{-10}$ & $3.9{\times}10^{-4}$ & 19 & 1 & 72 & 74 \\
B & 255 & $10^{4}$ & BoomerAMG & 8 & $6.7{\times}10^{-13}$ & $5.8{\times}10^{-9}$ & 5 & 38 & 45 & 84 \\
B & 255 & $10^{4}$ & M2 $w{=}\frac1{16}$ & 3 & $3.1{\times}10^{-11}$ & $1.1{\times}10^{-14}$ & 3 & 6 & 11 & 17 \\
B & 255 & $10^{4}$ & GMG & 7 & $2.8{\times}10^{-12}$ & $4.1{\times}10^{-10}$ & 5 & 21 & 35 & 56 \\
B & 255 & $10^{6}$ & 传统 LGF & 53 & $8.2{\times}10^{-10}$ & $6.9{\times}10^{-4}$ & 44 & 2 & 143 & 145 \\
B & 255 & $10^{6}$ & BoomerAMG & 8 & $6.7{\times}10^{-13}$ & $5.8{\times}10^{-9}$ & 5 & 39 & 45 & 84 \\
B & 255 & $10^{6}$ & M2 $w{=}\frac1{16}$ & 2 & $2.2{\times}10^{-7}$ & $1.1{\times}10^{-14}$ & 3 & 7 & 9 & 16 \\
B & 255 & $10^{6}$ & GMG & 7 & $3.0{\times}10^{-12}$ & $4.2{\times}10^{-10}$ & 5 & 20 & 31 & 52 \\
B & 255 & $10^{8}$ & 传统 LGF & 47 & $6.0{\times}10^{-7}$ & $5.6{\times}10^{-2}$ & 71 & 1 & 112 & 114 \\
B & 255 & $10^{8}$ & BoomerAMG & 8 & $6.7{\times}10^{-13}$ & $5.8{\times}10^{-9}$ & 5 & 39 & 44 & 83 \\
B & 255 & $10^{8}$ & M2 $w{=}\frac1{16}$ & 1 & $3.8{\times}10^{-4}$ & $1.1{\times}10^{-14}$ & 3 & 6 & 6 & 12 \\
B & 255 & $10^{8}$ & GMG & 7 & $3.0{\times}10^{-12}$ & $4.2{\times}10^{-10}$ & 5 & 20 & 29 & 49 \\
B & 511 & $10^{2}$ & 传统 LGF & 17 & $2.1{\times}10^{-12}$ & $3.6{\times}10^{-4}$ & 10 & 6 & 210 & 216 \\
B & 511 & $10^{2}$ & BoomerAMG & 8 & $6.7{\times}10^{-13}$ & $1.3{\times}10^{-8}$ & 6 & 155 & 183 & 338 \\
B & 511 & $10^{2}$ & M2 $w{=}\frac1{16}$ & 3 & $3.6{\times}10^{-11}$ & $1.2{\times}10^{-14}$ & 3 & 27 & 50 & 77 \\
B & 511 & $10^{2}$ & GMG & 7 & $2.6{\times}10^{-13}$ & $7.5{\times}10^{-11}$ & 5 & 81 & 135 & 216 \\
B & 511 & $10^{4}$ & 传统 LGF & 47 & $2.8{\times}10^{-12}$ & $3.9{\times}10^{-4}$ & 19 & 7 & 534 & 541 \\
B & 511 & $10^{4}$ & BoomerAMG & 8 & $2.8{\times}10^{-12}$ & $5.6{\times}10^{-9}$ & 5 & 154 & 180 & 333 \\
B & 511 & $10^{4}$ & M2 $w{=}\frac1{16}$ & 3 & $3.2{\times}10^{-11}$ & $1.1{\times}10^{-14}$ & 3 & 25 & 54 & 78 \\
B & 511 & $10^{4}$ & GMG & 7 & $2.3{\times}10^{-12}$ & $3.9{\times}10^{-11}$ & 5 & 77 & 121 & 197 \\
B & 511 & $10^{6}$ & 传统 LGF & 55 & $8.3{\times}10^{-10}$ & $4.3{\times}10^{-4}$ & 45 & 7 & 578 & 585 \\
B & 511 & $10^{6}$ & BoomerAMG & 8 & $2.9{\times}10^{-12}$ & $5.6{\times}10^{-9}$ & 5 & 146 & 177 & 323 \\
B & 511 & $10^{6}$ & M2 $w{=}\frac1{16}$ & 2 & $2.2{\times}10^{-7}$ & $1.1{\times}10^{-14}$ & 3 & 24 & 37 & 61 \\
B & 511 & $10^{6}$ & GMG & 7 & $2.6{\times}10^{-12}$ & $3.7{\times}10^{-11}$ & 5 & 82 & 137 & 219 \\
B & 511 & $10^{8}$ & 传统 LGF & 49 & $6.0{\times}10^{-7}$ & $3.2{\times}10^{-2}$ & 72 & 5 & 530 & 536 \\
B & 511 & $10^{8}$ & BoomerAMG & 8 & $2.9{\times}10^{-12}$ & $5.6{\times}10^{-9}$ & 5 & 150 & 174 & 323 \\
B & 511 & $10^{8}$ & M2 $w{=}\frac1{16}$ & 1 & $3.8{\times}10^{-4}$ & $1.1{\times}10^{-14}$ & 3 & 24 & 26 & 50 \\
B & 511 & $10^{8}$ & GMG & 7 & $2.6{\times}10^{-12}$ & $3.7{\times}10^{-11}$ & 5 & 76 & 122 & 198 \\
B & 1023 & $10^{2}$ & 传统 LGF & 19 & $2.0{\times}10^{-12}$ & $3.6{\times}10^{-4}$ & 10 & 23 & 1032 & 1055 \\
B & 1023 & $10^{2}$ & BoomerAMG & 9 & $3.4{\times}10^{-14}$ & $1.7{\times}10^{-8}$ & 6 & 569 & 748 & 1318 \\
B & 1023 & $10^{2}$ & M2 $w{=}\frac1{16}$ & 3 & $3.6{\times}10^{-11}$ & $1.2{\times}10^{-14}$ & 3 & 88 & 170 & 257 \\
B & 1023 & $10^{2}$ & GMG & 7 & $1.0{\times}10^{-13}$ & $4.2{\times}10^{-11}$ & 4 & 314 & 542 & 856 \\
B & 1023 & $10^{4}$ & 传统 LGF & 49 & $2.8{\times}10^{-12}$ & $3.9{\times}10^{-4}$ & 20 & 21 & 2663 & 2684 \\
B & 1023 & $10^{4}$ & BoomerAMG & 9 & $1.1{\times}10^{-13}$ & $8.6{\times}10^{-9}$ & 5 & 570 & 768 & 1338 \\
B & 1023 & $10^{4}$ & M2 $w{=}\frac1{16}$ & 3 & $3.2{\times}10^{-11}$ & $1.1{\times}10^{-14}$ & 3 & 91 & 267 & 358 \\
B & 1023 & $10^{4}$ & GMG & 7 & $1.8{\times}10^{-12}$ & $2.8{\times}10^{-10}$ & 5 & 323 & 555 & 878 \\
B & 1023 & $10^{6}$ & 传统 LGF & 58 & $8.3{\times}10^{-10}$ & $3.7{\times}10^{-4}$ & 46 & 26 & 3513 & 3539 \\
B & 1023 & $10^{6}$ & BoomerAMG & 9 & $1.1{\times}10^{-13}$ & $8.6{\times}10^{-9}$ & 5 & 596 & 787 & 1383 \\
B & 1023 & $10^{6}$ & M2 $w{=}\frac1{16}$ & 2 & $2.2{\times}10^{-7}$ & $1.1{\times}10^{-14}$ & 3 & 95 & 154 & 249 \\
B & 1023 & $10^{6}$ & GMG & 7 & $2.4{\times}10^{-12}$ & $3.4{\times}10^{-10}$ & 5 & 376 & 630 & 1007 \\
B & 1023 & $10^{8}$ & 传统 LGF & 52 & $6.0{\times}10^{-7}$ & $1.7{\times}10^{-2}$ & 72 & 25 & 3012 & 3038 \\
B & 1023 & $10^{8}$ & BoomerAMG & 9 & $1.1{\times}10^{-13}$ & $8.6{\times}10^{-9}$ & 5 & 584 & 767 & 1351 \\
B & 1023 & $10^{8}$ & M2 $w{=}\frac1{16}$ & 1 & $3.8{\times}10^{-4}$ & $1.1{\times}10^{-14}$ & 3 & 89 & 107 & 195 \\
B & 1023 & $10^{8}$ & GMG & 7 & $2.4{\times}10^{-12}$ & $3.4{\times}10^{-10}$ & 5 & 329 & 613 & 941 \\
C & 127 & $10^{2}$ & 传统 LGF & 1000$^\times$ & $1.6{\times}10^{-4}$ & $7.7{\times}10^{-1}$ & -- & 0 & 518 & 518 \\
C & 127 & $10^{2}$ & BoomerAMG & 9 & $3.3{\times}10^{-10}$ & $5.7{\times}10^{-6}$ & 8 & 12 & 14 & 26 \\
C & 127 & $10^{2}$ & M2 $w{=}\frac1{16}$ & 8 & $5.8{\times}10^{-8}$ & $3.8{\times}10^{-5}$ & 9 & 15 & 15 & 30 \\
C & 127 & $10^{2}$ & M2 $w{=}\frac1{32}$ & 11 & $3.3{\times}10^{-8}$ & $1.7{\times}10^{-4}$ & 12 & 10 & 16 & 26 \\
C & 127 & $10^{2}$ & M2 $w{=}\frac1{64}$ & 15 & $6.8{\times}10^{-9}$ & $5.9{\times}10^{-4}$ & 15 & 8 & 20 & 28 \\
C & 127 & $10^{2}$ & GMG & 58 & $1.5{\times}10^{-10}$ & $8.0{\times}10^{-1}$ & 56 & 6 & 63 & 69 \\
C & 127 & $10^{4}$ & 传统 LGF & 1000$^\times$ & $7.6{\times}10^{-2}$ & $8.5{\times}10^{-1}$ & -- & 0 & 554 & 554 \\
C & 127 & $10^{4}$ & BoomerAMG & 10 & $2.6{\times}10^{-10}$ & $2.0{\times}10^{-5}$ & 9 & 12 & 16 & 29 \\
C & 127 & $10^{4}$ & M2 $w{=}\frac1{16}$ & 8 & $1.4{\times}10^{-7}$ & $3.0{\times}10^{-5}$ & 9 & 15 & 15 & 29 \\
C & 127 & $10^{4}$ & M2 $w{=}\frac1{32}$ & 11 & $1.1{\times}10^{-7}$ & $3.9{\times}10^{-4}$ & 13 & 11 & 17 & 28 \\
C & 127 & $10^{4}$ & M2 $w{=}\frac1{64}$ & 15 & $2.9{\times}10^{-8}$ & $1.7{\times}10^{-3}$ & 16 & 9 & 21 & 30 \\
C & 127 & $10^{4}$ & GMG & 57 & $1.1{\times}10^{-9}$ & $7.8{\times}10^{-1}$ & 56 & 6 & 66 & 72 \\
C & 127 & $10^{6}$ & 传统 LGF & 1000$^\times$ & $1.5{\times}10^{-1}$ & $8.8{\times}10^{-1}$ & -- & 0 & 588 & 588 \\
C & 127 & $10^{6}$ & BoomerAMG & 9 & $2.6{\times}10^{-9}$ & $2.0{\times}10^{-5}$ & 9 & 12 & 14 & 27 \\
C & 127 & $10^{6}$ & M2 $w{=}\frac1{16}$ & 8 & $2.1{\times}10^{-7}$ & $2.4{\times}10^{-4}$ & 10 & 15 & 15 & 29 \\
C & 127 & $10^{6}$ & M2 $w{=}\frac1{32}$ & 11 & $2.8{\times}10^{-7}$ & $1.4{\times}10^{-3}$ & 13 & 11 & 15 & 26 \\
C & 127 & $10^{6}$ & M2 $w{=}\frac1{64}$ & 15 & $1.0{\times}10^{-7}$ & $6.6{\times}10^{-3}$ & 17 & 9 & 20 & 29 \\
C & 127 & $10^{6}$ & GMG & 49 & $2.5{\times}10^{-8}$ & $7.7{\times}10^{-1}$ & 56 & 6 & 58 & 64 \\
C & 127 & $10^{8}$ & 传统 LGF & 1000$^\times$ & $1.9{\times}10^{-1}$ & $8.9{\times}10^{-1}$ & -- & 0 & 531 & 532 \\
C & 127 & $10^{8}$ & BoomerAMG & 9 & $4.5{\times}10^{-10}$ & $2.1{\times}10^{-5}$ & 9 & 12 & 15 & 27 \\
C & 127 & $10^{8}$ & M2 $w{=}\frac1{16}$ & 8 & $2.5{\times}10^{-7}$ & $5.0{\times}10^{-4}$ & 10 & 15 & 16 & 31 \\
C & 127 & $10^{8}$ & M2 $w{=}\frac1{32}$ & 10 & $6.2{\times}10^{-7}$ & $4.4{\times}10^{-3}$ & 13 & 10 & 16 & 27 \\
C & 127 & $10^{8}$ & M2 $w{=}\frac1{64}$ & 14 & $4.0{\times}10^{-7}$ & $1.8{\times}10^{-2}$ & 18 & 9 & 20 & 29 \\
C & 127 & $10^{8}$ & GMG & 48 & $1.2{\times}10^{-6}$ & $7.4{\times}10^{-1}$ & 55 & 6 & 56 & 61 \\
C & 255 & $10^{2}$ & 传统 LGF & 1000$^\times$ & $2.8{\times}10^{-4}$ & $7.8{\times}10^{-1}$ & -- & 2 & 2286 & 2287 \\
C & 255 & $10^{2}$ & BoomerAMG & 9 & $7.7{\times}10^{-10}$ & $2.1{\times}10^{-5}$ & 8 & 45 & 54 & 99 \\
C & 255 & $10^{2}$ & M2 $w{=}\frac1{16}$ & 8 & $5.1{\times}10^{-8}$ & $3.5{\times}10^{-5}$ & 9 & 78 & 67 & 145 \\
C & 255 & $10^{2}$ & M2 $w{=}\frac1{32}$ & 11 & $2.1{\times}10^{-8}$ & $1.5{\times}10^{-4}$ & 12 & 44 & 68 & 112 \\
C & 255 & $10^{2}$ & M2 $w{=}\frac1{64}$ & 15 & $4.4{\times}10^{-9}$ & $5.1{\times}10^{-4}$ & 15 & 34 & 79 & 113 \\
C & 255 & $10^{2}$ & GMG & 68 & $7.0{\times}10^{-11}$ & $8.0{\times}10^{-1}$ & 64 & 19 & 272 & 292 \\
C & 255 & $10^{4}$ & 传统 LGF & 1000$^\times$ & $9.0{\times}10^{-2}$ & $8.5{\times}10^{-1}$ & -- & 1 & 2186 & 2188 \\
C & 255 & $10^{4}$ & BoomerAMG & 10 & $1.6{\times}10^{-10}$ & $8.2{\times}10^{-5}$ & 9 & 45 & 60 & 105 \\
C & 255 & $10^{4}$ & M2 $w{=}\frac1{16}$ & 8 & $8.4{\times}10^{-8}$ & $3.2{\times}10^{-5}$ & 9 & 75 & 70 & 145 \\
C & 255 & $10^{4}$ & M2 $w{=}\frac1{32}$ & 11 & $7.4{\times}10^{-8}$ & $3.4{\times}10^{-4}$ & 12 & 45 & 65 & 111 \\
C & 255 & $10^{4}$ & M2 $w{=}\frac1{64}$ & 14 & $5.5{\times}10^{-8}$ & $1.5{\times}10^{-3}$ & 16 & 34 & 73 & 106 \\
C & 255 & $10^{4}$ & GMG & 66 & $7.6{\times}10^{-10}$ & $7.8{\times}10^{-1}$ & 64 & 20 & 271 & 292 \\
C & 255 & $10^{6}$ & 传统 LGF & 1000$^\times$ & $1.5{\times}10^{-1}$ & $8.8{\times}10^{-1}$ & -- & 2 & 2086 & 2088 \\
C & 255 & $10^{6}$ & BoomerAMG & 10 & $8.6{\times}10^{-10}$ & $7.9{\times}10^{-5}$ & 9 & 48 & 62 & 110 \\
C & 255 & $10^{6}$ & M2 $w{=}\frac1{16}$ & 8 & $2.1{\times}10^{-7}$ & $1.6{\times}10^{-4}$ & 10 & 76 & 71 & 147 \\
C & 255 & $10^{6}$ & M2 $w{=}\frac1{32}$ & 11 & $1.7{\times}10^{-7}$ & $1.3{\times}10^{-3}$ & 13 & 44 & 69 & 112 \\
C & 255 & $10^{6}$ & M2 $w{=}\frac1{64}$ & 14 & $7.9{\times}10^{-8}$ & $5.6{\times}10^{-3}$ & 16 & 33 & 74 & 107 \\
C & 255 & $10^{6}$ & GMG & 65 & $5.6{\times}10^{-9}$ & $7.7{\times}10^{-1}$ & 65 & 19 & 269 & 289 \\
C & 255 & $10^{8}$ & 传统 LGF & 1000$^\times$ & $2.2{\times}10^{-1}$ & $8.9{\times}10^{-1}$ & -- & 1 & 2106 & 2107 \\
C & 255 & $10^{8}$ & BoomerAMG & 9 & $1.8{\times}10^{-9}$ & $6.1{\times}10^{-5}$ & 9 & 46 & 57 & 103 \\
C & 255 & $10^{8}$ & M2 $w{=}\frac1{16}$ & 8 & $2.7{\times}10^{-7}$ & $3.7{\times}10^{-4}$ & 10 & 76 & 73 & 150 \\
C & 255 & $10^{8}$ & M2 $w{=}\frac1{32}$ & 10 & $5.1{\times}10^{-7}$ & $3.6{\times}10^{-3}$ & 13 & 45 & 63 & 108 \\
C & 255 & $10^{8}$ & M2 $w{=}\frac1{64}$ & 14 & $2.2{\times}10^{-7}$ & $1.4{\times}10^{-2}$ & 17 & 35 & 75 & 110 \\
C & 255 & $10^{8}$ & GMG & 56 & $2.7{\times}10^{-7}$ & $7.5{\times}10^{-1}$ & 64 & 21 & 230 & 250 \\
C & 511 & $10^{2}$ & 传统 LGF & 1000$^\times$ & $3.9{\times}10^{-4}$ & $7.8{\times}10^{-1}$ & -- & 6 & 9435 & 9441 \\
C & 511 & $10^{2}$ & BoomerAMG & 10 & $7.5{\times}10^{-11}$ & $9.8{\times}10^{-6}$ & 9 & 169 & 228 & 397 \\
C & 511 & $10^{2}$ & M2 $w{=}\frac1{16}$ & 8 & $4.3{\times}10^{-8}$ & $3.1{\times}10^{-5}$ & 9 & 516 & 316 & 832 \\
C & 511 & $10^{2}$ & M2 $w{=}\frac1{32}$ & 11 & $1.6{\times}10^{-8}$ & $1.3{\times}10^{-4}$ & 12 & 230 & 270 & 499 \\
C & 511 & $10^{2}$ & M2 $w{=}\frac1{64}$ & 14 & $1.6{\times}10^{-8}$ & $4.5{\times}10^{-4}$ & 15 & 157 & 289 & 446 \\
C & 511 & $10^{2}$ & GMG & 87 & $1.4{\times}10^{-10}$ & $7.9{\times}10^{-1}$ & 81 & 78 & 1468 & 1546 \\
C & 511 & $10^{4}$ & 传统 LGF & 1000$^\times$ & $9.7{\times}10^{-2}$ & $8.6{\times}10^{-1}$ & -- & 5 & 8932 & 8937 \\
C & 511 & $10^{4}$ & BoomerAMG & 10 & $2.0{\times}10^{-10}$ & $9.9{\times}10^{-6}$ & 9 & 164 & 230 & 394 \\
C & 511 & $10^{4}$ & M2 $w{=}\frac1{16}$ & 8 & $6.8{\times}10^{-8}$ & $3.0{\times}10^{-5}$ & 9 & 509 & 318 & 826 \\
C & 511 & $10^{4}$ & M2 $w{=}\frac1{32}$ & 11 & $5.6{\times}10^{-8}$ & $3.0{\times}10^{-4}$ & 12 & 239 & 288 & 527 \\
C & 511 & $10^{4}$ & M2 $w{=}\frac1{64}$ & 14 & $3.4{\times}10^{-8}$ & $1.3{\times}10^{-3}$ & 16 & 160 & 296 & 456 \\
C & 511 & $10^{4}$ & GMG & 87 & $4.7{\times}10^{-10}$ & $7.7{\times}10^{-1}$ & 83 & 79 & 1408 & 1487 \\
C & 511 & $10^{6}$ & 传统 LGF & 1000$^\times$ & $1.6{\times}10^{-1}$ & $8.8{\times}10^{-1}$ & -- & 3 & 8111 & 8115 \\
C & 511 & $10^{6}$ & BoomerAMG & 10 & $2.2{\times}10^{-10}$ & $4.1{\times}10^{-5}$ & 9 & 182 & 250 & 432 \\
C & 511 & $10^{6}$ & M2 $w{=}\frac1{16}$ & 8 & $1.8{\times}10^{-7}$ & $1.3{\times}10^{-4}$ & 10 & 522 & 330 & 851 \\
C & 511 & $10^{6}$ & M2 $w{=}\frac1{32}$ & 10 & $3.4{\times}10^{-7}$ & $1.1{\times}10^{-3}$ & 13 & 240 & 266 & 505 \\
C & 511 & $10^{6}$ & M2 $w{=}\frac1{64}$ & 14 & $5.0{\times}10^{-8}$ & $4.8{\times}10^{-3}$ & 16 & 158 & 310 & 468 \\
C & 511 & $10^{6}$ & GMG & 85 & $1.8{\times}10^{-9}$ & $7.5{\times}10^{-1}$ & 84 & 78 & 1469 & 1547 \\
C & 511 & $10^{8}$ & 传统 LGF & 1000$^\times$ & $2.3{\times}10^{-1}$ & $8.9{\times}10^{-1}$ & -- & 4 & 8083 & 8087 \\
C & 511 & $10^{8}$ & BoomerAMG & 10 & $2.6{\times}10^{-10}$ & $1.7{\times}10^{-5}$ & 9 & 161 & 215 & 376 \\
C & 511 & $10^{8}$ & M2 $w{=}\frac1{16}$ & 8 & $2.3{\times}10^{-7}$ & $3.1{\times}10^{-4}$ & 10 & 502 & 305 & 806 \\
C & 511 & $10^{8}$ & M2 $w{=}\frac1{32}$ & 10 & $4.1{\times}10^{-7}$ & $3.1{\times}10^{-3}$ & 13 & 231 & 252 & 482 \\
C & 511 & $10^{8}$ & M2 $w{=}\frac1{64}$ & 13 & $7.9{\times}10^{-7}$ & $1.2{\times}10^{-2}$ & 16 & 155 & 261 & 416 \\
C & 511 & $10^{8}$ & GMG & 82 & $9.9{\times}10^{-9}$ & $7.3{\times}10^{-1}$ & 82 & 76 & 1299 & 1374 \\
C & 1023 & $10^{2}$ & 传统 LGF & 1000$^\times$ & $4.4{\times}10^{-4}$ & $7.8{\times}10^{-1}$ & -- & 32 & 43146 & 43177 \\
C & 1023 & $10^{2}$ & BoomerAMG & 11 & $5.0{\times}10^{-11}$ & $4.0{\times}10^{-5}$ & 9 & 653 & 975 & 1628 \\
C & 1023 & $10^{2}$ & M2 $w{=}\frac1{16}$ & 8 & $4.0{\times}10^{-8}$ & $3.0{\times}10^{-5}$ & 9 & 3445 & 1587 & 5032 \\
C & 1023 & $10^{2}$ & M2 $w{=}\frac1{32}$ & 11 & $1.4{\times}10^{-8}$ & $1.2{\times}10^{-4}$ & 12 & 1341 & 1393 & 2735 \\
C & 1023 & $10^{2}$ & M2 $w{=}\frac1{64}$ & 14 & $1.4{\times}10^{-8}$ & $4.2{\times}10^{-4}$ & 15 & 954 & 1386 & 2340 \\
C & 1023 & $10^{2}$ & GMG & 95 & $7.4{\times}10^{-11}$ & $7.9{\times}10^{-1}$ & 88 & 350 & 7236 & 7586 \\
C & 1023 & $10^{4}$ & 传统 LGF & 1000$^\times$ & $1.0{\times}10^{-1}$ & $8.6{\times}10^{-1}$ & -- & 23 & 38796 & 38819 \\
C & 1023 & $10^{4}$ & BoomerAMG & 11 & $8.8{\times}10^{-11}$ & $9.0{\times}10^{-5}$ & 9 & 648 & 952 & 1599 \\
C & 1023 & $10^{4}$ & M2 $w{=}\frac1{16}$ & 8 & $6.3{\times}10^{-8}$ & $2.8{\times}10^{-5}$ & 9 & 3438 & 1446 & 4884 \\
C & 1023 & $10^{4}$ & M2 $w{=}\frac1{32}$ & 11 & $5.1{\times}10^{-8}$ & $2.8{\times}10^{-4}$ & 12 & 1281 & 1271 & 2552 \\
C & 1023 & $10^{4}$ & M2 $w{=}\frac1{64}$ & 14 & $2.9{\times}10^{-8}$ & $1.3{\times}10^{-3}$ & 15 & 950 & 1406 & 2355 \\
C & 1023 & $10^{4}$ & GMG & 94 & $5.9{\times}10^{-10}$ & $7.7{\times}10^{-1}$ & 89 & 333 & 6587 & 6920 \\
C & 1023 & $10^{6}$ & 传统 LGF & 1000$^\times$ & $1.6{\times}10^{-1}$ & $8.8{\times}10^{-1}$ & -- & 19 & 39635 & 39654 \\
C & 1023 & $10^{6}$ & BoomerAMG & 10 & $4.7{\times}10^{-10}$ & $1.8{\times}10^{-5}$ & 9 & 630 & 858 & 1488 \\
C & 1023 & $10^{6}$ & M2 $w{=}\frac1{16}$ & 8 & $1.7{\times}10^{-7}$ & $1.3{\times}10^{-4}$ & 10 & 3481 & 1420 & 4901 \\
C & 1023 & $10^{6}$ & M2 $w{=}\frac1{32}$ & 10 & $3.1{\times}10^{-7}$ & $1.0{\times}10^{-3}$ & 13 & 1286 & 1139 & 2425 \\
C & 1023 & $10^{6}$ & M2 $w{=}\frac1{64}$ & 14 & $4.2{\times}10^{-8}$ & $4.6{\times}10^{-3}$ & 16 & 963 & 1401 & 2365 \\
C & 1023 & $10^{6}$ & GMG & 92 & $1.7{\times}10^{-9}$ & $7.5{\times}10^{-1}$ & 89 & 325 & 7338 & 7663 \\
C & 1023 & $10^{8}$ & 传统 LGF & 1000$^\times$ & $2.3{\times}10^{-1}$ & $8.9{\times}10^{-1}$ & -- & 23 & 38668 & 38691 \\
C & 1023 & $10^{8}$ & BoomerAMG & 10 & $9.2{\times}10^{-10}$ & $2.6{\times}10^{-5}$ & 9 & 651 & 878 & 1529 \\
C & 1023 & $10^{8}$ & M2 $w{=}\frac1{16}$ & 7 & $7.1{\times}10^{-6}$ & $3.0{\times}10^{-4}$ & 10 & 3454 & 1249 & 4703 \\
C & 1023 & $10^{8}$ & M2 $w{=}\frac1{32}$ & 10 & $3.8{\times}10^{-7}$ & $2.9{\times}10^{-3}$ & 13 & 1325 & 1263 & 2588 \\
C & 1023 & $10^{8}$ & M2 $w{=}\frac1{64}$ & 13 & $6.9{\times}10^{-7}$ & $1.1{\times}10^{-2}$ & 16 & 986 & 1277 & 2263 \\
C & 1023 & $10^{8}$ & GMG & 90 & $2.9{\times}10^{-9}$ & $7.3{\times}10^{-1}$ & 89 & 407 & 7796 & 8203 \\
D & 127 & $10^{2}$ & 传统 LGF & 1000$^\times$ & $3.3{\times}10^{-3}$ & $8.8{\times}10^{-1}$ & -- & 0 & 524 & 524 \\
D & 127 & $10^{2}$ & BoomerAMG & 9 & $2.6{\times}10^{-10}$ & $8.4{\times}10^{-6}$ & 8 & 20 & 19 & 40 \\
D & 127 & $10^{2}$ & M2 $w{=}\frac1{16}$ & 10 & $9.8{\times}10^{-9}$ & $1.3{\times}10^{-4}$ & 10 & 18 & 25 & 42 \\
D & 127 & $10^{2}$ & M2 $w{=}\frac1{32}$ & 15 & $1.4{\times}10^{-8}$ & $1.0{\times}10^{-3}$ & 16 & 13 & 26 & 39 \\
D & 127 & $10^{2}$ & M2 $w{=}\frac1{64}$ & 20 & $5.1{\times}10^{-8}$ & $4.4{\times}10^{-3}$ & 23 & 10 & 29 & 40 \\
D & 127 & $10^{2}$ & GMG & 80 & $9.3{\times}10^{-10}$ & $4.0{\times}10^{-1}$ & 77 & 6 & 87 & 93 \\
D & 127 & $10^{4}$ & 传统 LGF & 1000$^\times$ & $1.5{\times}10^{-1}$ & $9.6{\times}10^{-1}$ & -- & 1 & 611 & 612 \\
D & 127 & $10^{4}$ & BoomerAMG & 9 & $1.1{\times}10^{-10}$ & $9.8{\times}10^{-6}$ & 8 & 13 & 17 & 30 \\
D & 127 & $10^{4}$ & M2 $w{=}\frac1{16}$ & 9 & $8.5{\times}10^{-8}$ & $2.4{\times}10^{-4}$ & 11 & 17 & 21 & 38 \\
D & 127 & $10^{4}$ & M2 $w{=}\frac1{32}$ & 15 & $4.3{\times}10^{-8}$ & $2.3{\times}10^{-3}$ & 16 & 12 & 24 & 36 \\
D & 127 & $10^{4}$ & M2 $w{=}\frac1{64}$ & 21 & $5.1{\times}10^{-8}$ & $7.7{\times}10^{-3}$ & 23 & 11 & 31 & 43 \\
D & 127 & $10^{4}$ & GMG & 78 & $1.4{\times}10^{-9}$ & $3.9{\times}10^{-1}$ & 77 & 6 & 93 & 100 \\
D & 127 & $10^{6}$ & 传统 LGF & 1000$^\times$ & $2.6{\times}10^{-1}$ & $9.8{\times}10^{-1}$ & -- & 0 & 628 & 628 \\
D & 127 & $10^{6}$ & BoomerAMG & 9 & $1.3{\times}10^{-10}$ & $8.9{\times}10^{-5}$ & 8 & 13 & 17 & 30 \\
D & 127 & $10^{6}$ & M2 $w{=}\frac1{16}$ & 9 & $1.4{\times}10^{-7}$ & $4.0{\times}10^{-4}$ & 11 & 17 & 20 & 36 \\
D & 127 & $10^{6}$ & M2 $w{=}\frac1{32}$ & 13 & $2.4{\times}10^{-6}$ & $4.6{\times}10^{-3}$ & 17 & 12 & 23 & 35 \\
D & 127 & $10^{6}$ & M2 $w{=}\frac1{64}$ & 20 & $8.7{\times}10^{-7}$ & $1.4{\times}10^{-2}$ & 25 & 10 & 29 & 39 \\
D & 127 & $10^{6}$ & GMG & 76 & $1.8{\times}10^{-9}$ & $3.8{\times}10^{-1}$ & 75 & 6 & 81 & 86 \\
D & 127 & $10^{8}$ & 传统 LGF & 1000$^\times$ & $3.4{\times}10^{-1}$ & $9.8{\times}10^{-1}$ & -- & 0 & 627 & 627 \\
D & 127 & $10^{8}$ & BoomerAMG & 8 & $7.3{\times}10^{-10}$ & $1.4{\times}10^{-6}$ & 7 & 13 & 14 & 27 \\
D & 127 & $10^{8}$ & M2 $w{=}\frac1{16}$ & 9 & $7.6{\times}10^{-7}$ & $7.8{\times}10^{-4}$ & 11 & 17 & 20 & 36 \\
D & 127 & $10^{8}$ & M2 $w{=}\frac1{32}$ & 13 & $2.1{\times}10^{-6}$ & $1.4{\times}10^{-2}$ & 18 & 12 & 22 & 35 \\
D & 127 & $10^{8}$ & M2 $w{=}\frac1{64}$ & 19 & $1.8{\times}10^{-6}$ & $3.3{\times}10^{-2}$ & 26 & 10 & 30 & 40 \\
D & 127 & $10^{8}$ & GMG & 73 & $3.6{\times}10^{-8}$ & $3.6{\times}10^{-1}$ & 75 & 6 & 81 & 87 \\
D & 255 & $10^{2}$ & 传统 LGF & 1000$^\times$ & $4.9{\times}10^{-3}$ & $8.8{\times}10^{-1}$ & -- & 2 & 2267 & 2268 \\
D & 255 & $10^{2}$ & BoomerAMG & 10 & $3.5{\times}10^{-11}$ & $1.3{\times}10^{-5}$ & 8 & 48 & 65 & 112 \\
D & 255 & $10^{2}$ & M2 $w{=}\frac1{16}$ & 10 & $1.1{\times}10^{-8}$ & $1.1{\times}10^{-4}$ & 11 & 78 & 87 & 166 \\
D & 255 & $10^{2}$ & M2 $w{=}\frac1{32}$ & 14 & $5.0{\times}10^{-8}$ & $9.0{\times}10^{-4}$ & 16 & 46 & 90 & 136 \\
D & 255 & $10^{2}$ & M2 $w{=}\frac1{64}$ & 20 & $5.3{\times}10^{-8}$ & $3.2{\times}10^{-3}$ & 22 & 35 & 111 & 147 \\
D & 255 & $10^{2}$ & GMG & 93 & $8.5{\times}10^{-10}$ & $3.6{\times}10^{-1}$ & 89 & 23 & 390 & 412 \\
D & 255 & $10^{4}$ & 传统 LGF & 1000$^\times$ & $1.7{\times}10^{-1}$ & $9.6{\times}10^{-1}$ & -- & 2 & 2318 & 2320 \\
D & 255 & $10^{4}$ & BoomerAMG & 10 & $3.4{\times}10^{-11}$ & $4.9{\times}10^{-6}$ & 8 & 48 & 66 & 113 \\
D & 255 & $10^{4}$ & M2 $w{=}\frac1{16}$ & 9 & $1.1{\times}10^{-7}$ & $2.4{\times}10^{-4}$ & 11 & 79 & 83 & 162 \\
D & 255 & $10^{4}$ & M2 $w{=}\frac1{32}$ & 14 & $1.8{\times}10^{-7}$ & $2.3{\times}10^{-3}$ & 16 & 48 & 93 & 141 \\
D & 255 & $10^{4}$ & M2 $w{=}\frac1{64}$ & 21 & $4.8{\times}10^{-8}$ & $6.2{\times}10^{-3}$ & 23 & 36 & 117 & 153 \\
D & 255 & $10^{4}$ & GMG & 92 & $1.3{\times}10^{-9}$ & $3.4{\times}10^{-1}$ & 89 & 20 & 391 & 411 \\
D & 255 & $10^{6}$ & 传统 LGF & 1000$^\times$ & $2.9{\times}10^{-1}$ & $9.8{\times}10^{-1}$ & -- & 2 & 2192 & 2193 \\
D & 255 & $10^{6}$ & BoomerAMG & 10 & $1.5{\times}10^{-9}$ & $4.8{\times}10^{-5}$ & 9 & 48 & 61 & 110 \\
D & 255 & $10^{6}$ & M2 $w{=}\frac1{16}$ & 9 & $1.6{\times}10^{-7}$ & $4.3{\times}10^{-4}$ & 11 & 79 & 82 & 161 \\
D & 255 & $10^{6}$ & M2 $w{=}\frac1{32}$ & 13 & $1.7{\times}10^{-6}$ & $4.6{\times}10^{-3}$ & 17 & 48 & 88 & 136 \\
D & 255 & $10^{6}$ & M2 $w{=}\frac1{64}$ & 20 & $5.9{\times}10^{-7}$ & $1.2{\times}10^{-2}$ & 25 & 37 & 108 & 145 \\
D & 255 & $10^{6}$ & GMG & 89 & $5.0{\times}10^{-9}$ & $3.3{\times}10^{-1}$ & 88 & 22 & 379 & 401 \\
D & 255 & $10^{8}$ & 传统 LGF & 1000$^\times$ & $3.7{\times}10^{-1}$ & $9.8{\times}10^{-1}$ & -- & 1 & 2223 & 2224 \\
D & 255 & $10^{8}$ & BoomerAMG & 10 & $4.5{\times}10^{-10}$ & $2.7{\times}10^{-5}$ & 9 & 47 & 62 & 109 \\
D & 255 & $10^{8}$ & M2 $w{=}\frac1{16}$ & 9 & $7.7{\times}10^{-7}$ & $7.8{\times}10^{-4}$ & 11 & 78 & 82 & 160 \\
D & 255 & $10^{8}$ & M2 $w{=}\frac1{32}$ & 13 & $1.7{\times}10^{-6}$ & $1.3{\times}10^{-2}$ & 18 & 47 & 87 & 135 \\
D & 255 & $10^{8}$ & M2 $w{=}\frac1{64}$ & 19 & $1.5{\times}10^{-6}$ & $2.2{\times}10^{-2}$ & 26 & 36 & 107 & 143 \\
D & 255 & $10^{8}$ & GMG & 87 & $1.9{\times}10^{-8}$ & $3.1{\times}10^{-1}$ & 88 & 21 & 370 & 391 \\
D & 511 & $10^{2}$ & 传统 LGF & 1000$^\times$ & $5.6{\times}10^{-3}$ & $8.8{\times}10^{-1}$ & -- & 4 & 7503 & 7507 \\
D & 511 & $10^{2}$ & BoomerAMG & 10 & $2.7{\times}10^{-10}$ & $4.5{\times}10^{-5}$ & 9 & 175 & 235 & 409 \\
D & 511 & $10^{2}$ & M2 $w{=}\frac1{16}$ & 10 & $1.1{\times}10^{-8}$ & $1.0{\times}10^{-4}$ & 11 & 451 & 423 & 874 \\
D & 511 & $10^{2}$ & M2 $w{=}\frac1{32}$ & 14 & $3.2{\times}10^{-8}$ & $9.0{\times}10^{-4}$ & 15 & 220 & 365 & 585 \\
D & 511 & $10^{2}$ & M2 $w{=}\frac1{64}$ & 20 & $4.3{\times}10^{-8}$ & $2.7{\times}10^{-3}$ & 22 & 174 & 412 & 585 \\
D & 511 & $10^{2}$ & GMG & 104 & $4.3{\times}10^{-10}$ & $3.4{\times}10^{-1}$ & 100 & 79 & 1609 & 1687 \\
D & 511 & $10^{4}$ & 传统 LGF & 1000$^\times$ & $1.7{\times}10^{-1}$ & $9.6{\times}10^{-1}$ & -- & 5 & 7697 & 7703 \\
D & 511 & $10^{4}$ & BoomerAMG & 10 & $1.6{\times}10^{-10}$ & $1.7{\times}10^{-5}$ & 9 & 171 & 220 & 391 \\
D & 511 & $10^{4}$ & M2 $w{=}\frac1{16}$ & 9 & $1.2{\times}10^{-7}$ & $2.3{\times}10^{-4}$ & 11 & 444 & 372 & 816 \\
D & 511 & $10^{4}$ & M2 $w{=}\frac1{32}$ & 14 & $1.4{\times}10^{-7}$ & $2.2{\times}10^{-3}$ & 16 & 223 & 383 & 606 \\
D & 511 & $10^{4}$ & M2 $w{=}\frac1{64}$ & 20 & $1.2{\times}10^{-7}$ & $5.3{\times}10^{-3}$ & 23 & 155 & 426 & 581 \\
D & 511 & $10^{4}$ & GMG & 103 & $6.1{\times}10^{-10}$ & $3.2{\times}10^{-1}$ & 100 & 78 & 1670 & 1747 \\
D & 511 & $10^{6}$ & 传统 LGF & 1000$^\times$ & $3.0{\times}10^{-1}$ & $9.8{\times}10^{-1}$ & -- & 6 & 8303 & 8309 \\
D & 511 & $10^{6}$ & BoomerAMG & 10 & $2.0{\times}10^{-10}$ & $1.7{\times}10^{-5}$ & 9 & 171 & 233 & 404 \\
D & 511 & $10^{6}$ & M2 $w{=}\frac1{16}$ & 9 & $1.3{\times}10^{-7}$ & $4.0{\times}10^{-4}$ & 11 & 454 & 372 & 826 \\
D & 511 & $10^{6}$ & M2 $w{=}\frac1{32}$ & 13 & $1.7{\times}10^{-6}$ & $4.3{\times}10^{-3}$ & 17 & 231 & 360 & 591 \\
D & 511 & $10^{6}$ & M2 $w{=}\frac1{64}$ & 19 & $1.0{\times}10^{-6}$ & $1.1{\times}10^{-2}$ & 24 & 155 & 417 & 572 \\
D & 511 & $10^{6}$ & GMG & 100 & $3.8{\times}10^{-9}$ & $3.1{\times}10^{-1}$ & 99 & 80 & 1661 & 1741 \\
D & 511 & $10^{8}$ & 传统 LGF & 1000$^\times$ & $3.9{\times}10^{-1}$ & $9.8{\times}10^{-1}$ & -- & 4 & 8006 & 8010 \\
D & 511 & $10^{8}$ & BoomerAMG & 10 & $1.5{\times}10^{-8}$ & $8.5{\times}10^{-5}$ & 11 & 194 & 250 & 444 \\
D & 511 & $10^{8}$ & M2 $w{=}\frac1{16}$ & 9 & $6.0{\times}10^{-7}$ & $7.2{\times}10^{-4}$ & 11 & 482 & 409 & 891 \\
D & 511 & $10^{8}$ & M2 $w{=}\frac1{32}$ & 13 & $2.0{\times}10^{-6}$ & $1.2{\times}10^{-2}$ & 18 & 248 & 392 & 640 \\
D & 511 & $10^{8}$ & M2 $w{=}\frac1{64}$ & 19 & $1.3{\times}10^{-6}$ & $2.0{\times}10^{-2}$ & 25 & 165 & 432 & 597 \\
D & 511 & $10^{8}$ & GMG & 98 & $8.2{\times}10^{-9}$ & $2.9{\times}10^{-1}$ & 98 & 79 & 1724 & 1802 \\
D & 1023 & $10^{2}$ & 传统 LGF & 1000$^\times$ & $6.6{\times}10^{-3}$ & $8.8{\times}10^{-1}$ & -- & 21 & 39572 & 39593 \\
D & 1023 & $10^{2}$ & BoomerAMG & 10 & $1.4{\times}10^{-10}$ & $2.1{\times}10^{-5}$ & 9 & 632 & 840 & 1472 \\
D & 1023 & $10^{2}$ & M2 $w{=}\frac1{16}$ & 10 & $1.1{\times}10^{-8}$ & $9.8{\times}10^{-5}$ & 11 & 3149 & 1837 & 4986 \\
D & 1023 & $10^{2}$ & M2 $w{=}\frac1{32}$ & 14 & $3.3{\times}10^{-8}$ & $8.6{\times}10^{-4}$ & 15 & 1295 & 1637 & 2932 \\
D & 1023 & $10^{2}$ & M2 $w{=}\frac1{64}$ & 20 & $3.9{\times}10^{-8}$ & $2.5{\times}10^{-3}$ & 22 & 822 & 1851 & 2673 \\
D & 1023 & $10^{2}$ & GMG & 114 & $1.2{\times}10^{-9}$ & $3.4{\times}10^{-1}$ & 108 & 386 & 9984 & 10370 \\
D & 1023 & $10^{4}$ & 传统 LGF & 1000$^\times$ & $1.8{\times}10^{-1}$ & $9.6{\times}10^{-1}$ & -- & 22 & 38472 & 38494 \\
D & 1023 & $10^{4}$ & BoomerAMG & 11 & $4.2{\times}10^{-11}$ & $2.5{\times}10^{-5}$ & 9 & 644 & 961 & 1605 \\
D & 1023 & $10^{4}$ & M2 $w{=}\frac1{16}$ & 9 & $1.2{\times}10^{-7}$ & $2.2{\times}10^{-4}$ & 11 & 3183 & 1763 & 4946 \\
D & 1023 & $10^{4}$ & M2 $w{=}\frac1{32}$ & 14 & $1.4{\times}10^{-7}$ & $2.1{\times}10^{-3}$ & 16 & 1383 & 1808 & 3191 \\
D & 1023 & $10^{4}$ & M2 $w{=}\frac1{64}$ & 20 & $1.0{\times}10^{-7}$ & $5.0{\times}10^{-3}$ & 23 & 861 & 1978 & 2839 \\
D & 1023 & $10^{4}$ & GMG & 112 & $4.5{\times}10^{-10}$ & $3.2{\times}10^{-1}$ & 107 & 543 & 11165 & 11708 \\
D & 1023 & $10^{6}$ & 传统 LGF & 1000$^\times$ & $3.1{\times}10^{-1}$ & $9.8{\times}10^{-1}$ & -- & 29 & 42240 & 42269 \\
D & 1023 & $10^{6}$ & BoomerAMG & 10 & $2.5{\times}10^{-10}$ & $1.8{\times}10^{-5}$ & 9 & 682 & 898 & 1580 \\
D & 1023 & $10^{6}$ & M2 $w{=}\frac1{16}$ & 9 & $1.4{\times}10^{-7}$ & $3.9{\times}10^{-4}$ & 11 & 3422 & 1845 & 5267 \\
D & 1023 & $10^{6}$ & M2 $w{=}\frac1{32}$ & 13 & $1.6{\times}10^{-6}$ & $4.2{\times}10^{-3}$ & 17 & 1459 & 1761 & 3219 \\
D & 1023 & $10^{6}$ & M2 $w{=}\frac1{64}$ & 19 & $9.1{\times}10^{-7}$ & $1.0{\times}10^{-2}$ & 24 & 850 & 1924 & 2774 \\
D & 1023 & $10^{6}$ & GMG & 111 & $7.4{\times}10^{-10}$ & $3.0{\times}10^{-1}$ & 108 & 357 & 9403 & 9760 \\
D & 1023 & $10^{8}$ & 传统 LGF & 1000$^\times$ & $4.0{\times}10^{-1}$ & $9.8{\times}10^{-1}$ & -- & 19 & 37898 & 37917 \\
D & 1023 & $10^{8}$ & BoomerAMG & 10 & $1.5{\times}10^{-10}$ & $4.0{\times}10^{-5}$ & 9 & 683 & 939 & 1622 \\
D & 1023 & $10^{8}$ & M2 $w{=}\frac1{16}$ & 8 & $3.0{\times}10^{-6}$ & $7.2{\times}10^{-4}$ & 11 & 3356 & 1660 & 5016 \\
D & 1023 & $10^{8}$ & M2 $w{=}\frac1{32}$ & 13 & $1.9{\times}10^{-6}$ & $1.1{\times}10^{-2}$ & 18 & 1403 & 1744 & 3147 \\
D & 1023 & $10^{8}$ & M2 $w{=}\frac1{64}$ & 19 & $1.2{\times}10^{-6}$ & $1.9{\times}10^{-2}$ & 25 & 866 & 1855 & 2721 \\
D & 1023 & $10^{8}$ & GMG & 108 & $4.9{\times}10^{-9}$ & $2.8{\times}10^{-1}$ & 107 & 351 & 8977 & 9328 \\
\bottomrule
\end{longtable}
\endgroup
\IfFileExists{appendix_3d.tex}{\begingroup\footnotesize\setlength{\tabcolsep}{3pt}
\begin{longtable}{lrrlrcccrrr}
\caption{三维 全部结果。err：按残差规则停止时相对 $u_h$ 的代数误差；$k_8$：误差降到 $10^{-8}$ 所需迭代数；时间单位 ms（单线程）。}\\
\toprule 算例 & $n$ & $\chi$ & 方法 & 迭代 & err & 第 5 步误差 & $k_8$ & setup & solve & 合计 \\ \midrule
\endfirsthead \toprule 算例 & $n$ & $\chi$ & 方法 & 迭代 & err & 第 5 步误差 & $k_8$ & setup & solve & 合计 \\ \midrule \endhead
A & 31 & $10^{2}$ & 传统 LGF & 68 & $8.2{\times}10^{-12}$ & $1.5{\times}10^{-1}$ & 54 & 0 & 90 & 90 \\
A & 31 & $10^{2}$ & GMG & 9 & $4.2{\times}10^{-12}$ & $8.9{\times}10^{-8}$ & 6 & 10 & 20 & 29 \\
A & 31 & $10^{2}$ & BoomerAMG & 7 & $1.8{\times}10^{-13}$ & $6.7{\times}10^{-10}$ & 5 & 40 & 33 & 73 \\
A & 31 & $10^{2}$ & 方法二 & 3 & $1.2{\times}10^{-12}$ & $1.4{\times}10^{-16}$ & 3 & 0 & 6 & 6 \\
A & 31 & $10^{4}$ & 传统 LGF & 324 & $3.2{\times}10^{-10}$ & $5.0{\times}10^{-1}$ & 244 & 0 & 436 & 436 \\
A & 31 & $10^{4}$ & GMG & 9 & $3.9{\times}10^{-11}$ & $3.7{\times}10^{-7}$ & 7 & 9 & 20 & 29 \\
A & 31 & $10^{4}$ & BoomerAMG & 6 & $1.7{\times}10^{-11}$ & $8.4{\times}10^{-10}$ & 5 & 58 & 42 & 100 \\
A & 31 & $10^{4}$ & 方法二 & 2 & $3.3{\times}10^{-9}$ & $1.0{\times}10^{-16}$ & 2 & 1 & 5 & 5 \\
A & 31 & $10^{6}$ & 传统 LGF & 500$^\times$ & $9.9{\times}10^{-4}$ & $6.6{\times}10^{-1}$ & -- & 0 & 675 & 675 \\
A & 31 & $10^{6}$ & GMG & 9 & $2.5{\times}10^{-10}$ & $1.0{\times}10^{-6}$ & 8 & 10 & 20 & 31 \\
A & 31 & $10^{6}$ & BoomerAMG & 6 & $1.1{\times}10^{-11}$ & $1.0{\times}10^{-9}$ & 5 & 60 & 43 & 104 \\
A & 31 & $10^{6}$ & 方法二 & 2 & $7.5{\times}10^{-10}$ & $3.9{\times}10^{-14}$ & 2 & 1 & 4 & 5 \\
A & 31 & $10^{8}$ & 传统 LGF & 500$^\times$ & $2.4{\times}10^{-2}$ & $7.4{\times}10^{-1}$ & -- & 0 & 745 & 745 \\
A & 31 & $10^{8}$ & GMG & 8 & $6.1{\times}10^{-9}$ & $2.0{\times}10^{-6}$ & 8 & 10 & 22 & 33 \\
A & 31 & $10^{8}$ & BoomerAMG & 5 & $1.7{\times}10^{-10}$ & $1.7{\times}10^{-10}$ & 5 & 62 & 38 & 100 \\
A & 31 & $10^{8}$ & 方法二 & 2 & $2.6{\times}10^{-10}$ & $1.1{\times}10^{-14}$ & 2 & 1 & 4 & 5 \\
A & 63 & $10^{2}$ & 传统 LGF & 97 & $1.9{\times}10^{-12}$ & $1.6{\times}10^{-1}$ & 65 & 1 & 1145 & 1146 \\
A & 63 & $10^{2}$ & GMG & 9 & $2.5{\times}10^{-12}$ & $6.1{\times}10^{-8}$ & 6 & 77 & 165 & 241 \\
A & 63 & $10^{2}$ & BoomerAMG & 7 & $4.3{\times}10^{-12}$ & $5.8{\times}10^{-9}$ & 5 & 354 & 317 & 671 \\
A & 63 & $10^{2}$ & 方法二 & 3 & $1.1{\times}10^{-12}$ & $1.3{\times}10^{-16}$ & 3 & 4 & 54 & 59 \\
A & 63 & $10^{4}$ & 传统 LGF & 500$^\times$ & $2.5{\times}10^{-9}$ & $5.2{\times}10^{-1}$ & -- & 1 & 6007 & 6008 \\
A & 63 & $10^{4}$ & GMG & 9 & $1.0{\times}10^{-11}$ & $2.1{\times}10^{-7}$ & 7 & 90 & 198 & 288 \\
A & 63 & $10^{4}$ & BoomerAMG & 7 & $7.1{\times}10^{-12}$ & $1.5{\times}10^{-8}$ & 6 & 355 & 315 & 670 \\
A & 63 & $10^{4}$ & 方法二 & 2 & $3.3{\times}10^{-9}$ & $1.0{\times}10^{-16}$ & 2 & 4 & 39 & 43 \\
A & 63 & $10^{6}$ & 传统 LGF & 500$^\times$ & $4.6{\times}10^{-3}$ & $7.0{\times}10^{-1}$ & -- & 1 & 6447 & 6448 \\
A & 63 & $10^{6}$ & GMG & 9 & $3.1{\times}10^{-11}$ & $5.6{\times}10^{-7}$ & 7 & 85 & 181 & 266 \\
A & 63 & $10^{6}$ & BoomerAMG & 7 & $1.5{\times}10^{-10}$ & $1.6{\times}10^{-7}$ & 6 & 685 & 543 & 1228 \\
A & 63 & $10^{6}$ & 方法二 & 2 & $7.4{\times}10^{-10}$ & $3.8{\times}10^{-14}$ & 2 & 5 & 45 & 50 \\
A & 63 & $10^{8}$ & 传统 LGF & 500$^\times$ & $5.8{\times}10^{-2}$ & $7.8{\times}10^{-1}$ & -- & 1 & 6523 & 6524 \\
A & 63 & $10^{8}$ & GMG & 9 & $7.1{\times}10^{-11}$ & $9.2{\times}10^{-7}$ & 7 & 84 & 173 & 257 \\
A & 63 & $10^{8}$ & BoomerAMG & 7 & $6.5{\times}10^{-11}$ & $7.3{\times}10^{-8}$ & 6 & 673 & 559 & 1232 \\
A & 63 & $10^{8}$ & 方法二 & 2 & $2.6{\times}10^{-10}$ & $1.1{\times}10^{-14}$ & 2 & 4 & 37 & 41 \\
A & 127 & $10^{2}$ & 传统 LGF & 111 & $6.5{\times}10^{-13}$ & $1.6{\times}10^{-1}$ & 70 & 4 & 9966 & 9970 \\
A & 127 & $10^{2}$ & GMG & 9 & $1.8{\times}10^{-12}$ & $3.1{\times}10^{-8}$ & 6 & 773 & 1724 & 2496 \\
A & 127 & $10^{2}$ & BoomerAMG & 8 & $1.2{\times}10^{-12}$ & $2.2{\times}10^{-8}$ & 6 & 3873 & 4194 & 8067 \\
A & 127 & $10^{2}$ & 方法二 & 3 & $1.1{\times}10^{-12}$ & $1.3{\times}10^{-16}$ & 3 & 35 & 407 & 442 \\
A & 127 & $10^{4}$ & 传统 LGF & 500$^\times$ & $5.6{\times}10^{-8}$ & $5.3{\times}10^{-1}$ & -- & 4 & 47938 & 47942 \\
A & 127 & $10^{4}$ & GMG & 9 & $5.6{\times}10^{-12}$ & $1.4{\times}10^{-7}$ & 7 & 755 & 1482 & 2237 \\
A & 127 & $10^{4}$ & BoomerAMG & 8 & $1.3{\times}10^{-12}$ & $7.0{\times}10^{-8}$ & 6 & 3406 & 3646 & 7052 \\
A & 127 & $10^{4}$ & 方法二 & 2 & $3.3{\times}10^{-9}$ & $1.0{\times}10^{-16}$ & 2 & 34 & 303 & 337 \\
A & 127 & $10^{6}$ & 传统 LGF & 500$^\times$ & $9.1{\times}10^{-3}$ & $7.1{\times}10^{-1}$ & -- & 4 & 47288 & 47292 \\
A & 127 & $10^{6}$ & GMG & 9 & $1.6{\times}10^{-11}$ & $3.9{\times}10^{-7}$ & 7 & 1022 & 1830 & 2852 \\
A & 127 & $10^{6}$ & BoomerAMG & 7 & $1.5{\times}10^{-10}$ & $1.3{\times}10^{-7}$ & 6 & 4325 & 3671 & 7996 \\
A & 127 & $10^{6}$ & 方法二 & 2 & $7.4{\times}10^{-10}$ & $3.8{\times}10^{-14}$ & 2 & 41 & 308 & 350 \\
A & 127 & $10^{8}$ & 传统 LGF & 500$^\times$ & $8.4{\times}10^{-2}$ & $8.0{\times}10^{-1}$ & -- & 4 & 47956 & 47960 \\
A & 127 & $10^{8}$ & GMG & 9 & $3.2{\times}10^{-11}$ & $6.1{\times}10^{-7}$ & 7 & 945 & 1674 & 2619 \\
A & 127 & $10^{8}$ & BoomerAMG & 9 & $5.2{\times}10^{-11}$ & $6.9{\times}10^{-6}$ & 8 & 4734 & 4125 & 8859 \\
A & 127 & $10^{8}$ & 方法二 & 2 & $2.6{\times}10^{-10}$ & $1.1{\times}10^{-14}$ & 2 & 34 & 301 & 335 \\
B & 31 & $10^{2}$ & 传统 LGF & 20 & $2.1{\times}10^{-12}$ & $5.0{\times}10^{-4}$ & 12 & 0 & 31 & 31 \\
B & 31 & $10^{2}$ & GMG & 9 & $3.1{\times}10^{-12}$ & $3.9{\times}10^{-8}$ & 6 & 10 & 22 & 32 \\
B & 31 & $10^{2}$ & BoomerAMG & 8 & $2.3{\times}10^{-13}$ & $5.9{\times}10^{-8}$ & 6 & 70 & 55 & 124 \\
B & 31 & $10^{2}$ & 方法二 & 3 & $1.2{\times}10^{-11}$ & $3.2{\times}10^{-15}$ & 3 & 0 & 6 & 7 \\
B & 31 & $10^{4}$ & 传统 LGF & 29 & $8.3{\times}10^{-10}$ & $5.4{\times}10^{-4}$ & 19 & 0 & 42 & 42 \\
B & 31 & $10^{4}$ & GMG & 9 & $3.3{\times}10^{-12}$ & $4.2{\times}10^{-8}$ & 6 & 10 & 21 & 31 \\
B & 31 & $10^{4}$ & BoomerAMG & 7 & $3.7{\times}10^{-10}$ & $8.4{\times}10^{-8}$ & 6 & 71 & 52 & 123 \\
B & 31 & $10^{4}$ & 方法二 & 3 & $1.1{\times}10^{-11}$ & $3.1{\times}10^{-15}$ & 3 & 0 & 6 & 6 \\
B & 31 & $10^{6}$ & 传统 LGF & 31 & $7.3{\times}10^{-7}$ & $2.1{\times}10^{-3}$ & 49 & 0 & 50 & 50 \\
B & 31 & $10^{6}$ & GMG & 9 & $3.3{\times}10^{-12}$ & $4.2{\times}10^{-8}$ & 6 & 11 & 23 & 34 \\
B & 31 & $10^{6}$ & BoomerAMG & 7 & $2.7{\times}10^{-10}$ & $1.6{\times}10^{-7}$ & 6 & 69 & 50 & 119 \\
B & 31 & $10^{6}$ & 方法二 & 2 & $9.7{\times}10^{-8}$ & $3.1{\times}10^{-15}$ & 3 & 0 & 4 & 5 \\
B & 31 & $10^{8}$ & 传统 LGF & 48 & $7.3{\times}10^{-7}$ & $1.3{\times}10^{-1}$ & 80 & 0 & 69 & 69 \\
B & 31 & $10^{8}$ & GMG & 9 & $3.3{\times}10^{-12}$ & $4.2{\times}10^{-8}$ & 6 & 10 & 21 & 31 \\
B & 31 & $10^{8}$ & BoomerAMG & 7 & $4.1{\times}10^{-10}$ & $1.0{\times}10^{-7}$ & 6 & 70 & 49 & 119 \\
B & 31 & $10^{8}$ & 方法二 & 1 & $2.0{\times}10^{-4}$ & $3.1{\times}10^{-15}$ & 3 & 1 & 3 & 4 \\
B & 63 & $10^{2}$ & 传统 LGF & 19 & $2.2{\times}10^{-12}$ & $5.1{\times}10^{-4}$ & 12 & 1 & 279 & 280 \\
B & 63 & $10^{2}$ & GMG & 9 & $8.3{\times}10^{-13}$ & $1.7{\times}10^{-8}$ & 6 & 85 & 185 & 270 \\
B & 63 & $10^{2}$ & BoomerAMG & 9 & $3.6{\times}10^{-12}$ & $7.7{\times}10^{-7}$ & 7 & 662 & 559 & 1221 \\
B & 63 & $10^{2}$ & 方法二 & 3 & $1.2{\times}10^{-11}$ & $3.2{\times}10^{-15}$ & 3 & 5 & 57 & 62 \\
B & 63 & $10^{4}$ & 传统 LGF & 31 & $8.6{\times}10^{-10}$ & $5.4{\times}10^{-4}$ & 19 & 1 & 421 & 421 \\
B & 63 & $10^{4}$ & GMG & 9 & $8.9{\times}10^{-13}$ & $1.8{\times}10^{-8}$ & 6 & 85 & 178 & 263 \\
B & 63 & $10^{4}$ & BoomerAMG & 9 & $6.6{\times}10^{-11}$ & $7.6{\times}10^{-7}$ & 7 & 668 & 568 & 1235 \\
B & 63 & $10^{4}$ & 方法二 & 3 & $1.1{\times}10^{-11}$ & $3.0{\times}10^{-15}$ & 3 & 4 & 53 & 58 \\
B & 63 & $10^{6}$ & 传统 LGF & 59 & $8.7{\times}10^{-10}$ & $1.1{\times}10^{-3}$ & 51 & 1 & 760 & 760 \\
B & 63 & $10^{6}$ & GMG & 9 & $8.9{\times}10^{-13}$ & $1.8{\times}10^{-8}$ & 6 & 83 & 197 & 280 \\
B & 63 & $10^{6}$ & BoomerAMG & 9 & $1.0{\times}10^{-10}$ & $1.2{\times}10^{-6}$ & 7 & 654 & 559 & 1214 \\
B & 63 & $10^{6}$ & 方法二 & 2 & $9.6{\times}10^{-8}$ & $3.0{\times}10^{-15}$ & 3 & 4 & 37 & 41 \\
B & 63 & $10^{8}$ & 传统 LGF & 53 & $7.5{\times}10^{-7}$ & $8.2{\times}10^{-2}$ & 85 & 1 & 697 & 698 \\
B & 63 & $10^{8}$ & GMG & 9 & $8.9{\times}10^{-13}$ & $1.8{\times}10^{-8}$ & 6 & 86 & 179 & 265 \\
B & 63 & $10^{8}$ & BoomerAMG & 9 & $1.6{\times}10^{-10}$ & $1.2{\times}10^{-6}$ & 8 & 649 & 559 & 1208 \\
B & 63 & $10^{8}$ & 方法二 & 1 & $2.0{\times}10^{-4}$ & $3.0{\times}10^{-15}$ & 3 & 4 & 25 & 29 \\
B & 127 & $10^{2}$ & 传统 LGF & 19 & $2.0{\times}10^{-12}$ & $5.1{\times}10^{-4}$ & 12 & 5 & 1853 & 1858 \\
B & 127 & $10^{2}$ & GMG & 9 & $2.9{\times}10^{-13}$ & $6.9{\times}10^{-9}$ & 5 & 746 & 1470 & 2216 \\
B & 127 & $10^{2}$ & BoomerAMG & 11 & $9.6{\times}10^{-13}$ & $7.0{\times}10^{-6}$ & 8 & 6157 & 7380 & 13537 \\
B & 127 & $10^{2}$ & 方法二 & 3 & $1.2{\times}10^{-11}$ & $3.1{\times}10^{-15}$ & 3 & 34 & 395 & 429 \\
B & 127 & $10^{4}$ & 传统 LGF & 33 & $8.6{\times}10^{-10}$ & $5.5{\times}10^{-4}$ & 19 & 6 & 3448 & 3454 \\
B & 127 & $10^{4}$ & GMG & 9 & $3.4{\times}10^{-13}$ & $7.4{\times}10^{-9}$ & 5 & 791 & 1463 & 2254 \\
B & 127 & $10^{4}$ & BoomerAMG & 11 & $1.5{\times}10^{-11}$ & $9.3{\times}10^{-6}$ & 9 & 5728 & 7164 & 12892 \\
B & 127 & $10^{4}$ & 方法二 & 3 & $1.1{\times}10^{-11}$ & $3.0{\times}10^{-15}$ & 3 & 39 & 427 & 466 \\
B & 127 & $10^{6}$ & 传统 LGF & 62 & $8.7{\times}10^{-10}$ & $6.5{\times}10^{-4}$ & 53 & 7 & 5818 & 5825 \\
B & 127 & $10^{6}$ & GMG & 9 & $3.4{\times}10^{-13}$ & $7.4{\times}10^{-9}$ & 5 & 755 & 1466 & 2221 \\
B & 127 & $10^{6}$ & BoomerAMG & 10 & $3.9{\times}10^{-10}$ & $6.9{\times}10^{-6}$ & 9 & 6041 & 6186 & 12227 \\
B & 127 & $10^{6}$ & 方法二 & 2 & $9.6{\times}10^{-8}$ & $3.0{\times}10^{-15}$ & 3 & 34 & 296 & 329 \\
B & 127 & $10^{8}$ & 传统 LGF & 56 & $7.5{\times}10^{-7}$ & $4.8{\times}10^{-2}$ & 87 & 5 & 5131 & 5136 \\
B & 127 & $10^{8}$ & GMG & 9 & $3.4{\times}10^{-13}$ & $7.4{\times}10^{-9}$ & 5 & 767 & 1509 & 2277 \\
B & 127 & $10^{8}$ & BoomerAMG & 11 & $4.8{\times}10^{-11}$ & $6.2{\times}10^{-6}$ & 9 & 5998 & 6798 & 12796 \\
B & 127 & $10^{8}$ & 方法二 & 1 & $2.0{\times}10^{-4}$ & $3.0{\times}10^{-15}$ & 3 & 38 & 191 & 229 \\
C & 31 & $10^{2}$ & 传统 LGF & 500$^\times$ & $8.5{\times}10^{-3}$ & $8.0{\times}10^{-1}$ & -- & 0 & 712 & 712 \\
C & 31 & $10^{2}$ & GMG & 62 & $1.4{\times}10^{-8}$ & $6.0{\times}10^{-1}$ & 67 & 11 & 149 & 160 \\
C & 31 & $10^{2}$ & BoomerAMG & 11 & $2.2{\times}10^{-10}$ & $7.8{\times}10^{-5}$ & 10 & 51 & 63 & 114 \\
C & 31 & $10^{2}$ & M2 $w{=}\frac1{32}$ & 12 & $1.5{\times}10^{-8}$ & $2.9{\times}10^{-4}$ & 13 & 14 & 34 & 49 \\
C & 31 & $10^{2}$ & M2 $w{=}\frac1{64}$ & 12 & $1.5{\times}10^{-8}$ & $2.9{\times}10^{-4}$ & 13 & 14 & 34 & 48 \\
C & 31 & $10^{4}$ & 传统 LGF & 500$^\times$ & $1.8{\times}10^{-1}$ & $8.7{\times}10^{-1}$ & -- & 0 & 684 & 684 \\
C & 31 & $10^{4}$ & GMG & 61 & $5.1{\times}10^{-8}$ & $5.6{\times}10^{-1}$ & 66 & 10 & 139 & 149 \\
C & 31 & $10^{4}$ & BoomerAMG & 11 & $4.8{\times}10^{-10}$ & $7.8{\times}10^{-5}$ & 10 & 53 & 66 & 118 \\
C & 31 & $10^{4}$ & M2 $w{=}\frac1{32}$ & 12 & $5.1{\times}10^{-8}$ & $1.0{\times}10^{-3}$ & 13 & 14 & 34 & 48 \\
C & 31 & $10^{4}$ & M2 $w{=}\frac1{64}$ & 12 & $5.1{\times}10^{-8}$ & $1.0{\times}10^{-3}$ & 13 & 15 & 36 & 50 \\
C & 31 & $10^{6}$ & 传统 LGF & 500$^\times$ & $2.5{\times}10^{-1}$ & $9.0{\times}10^{-1}$ & -- & 0 & 746 & 746 \\
C & 31 & $10^{6}$ & GMG & 60 & $5.9{\times}10^{-8}$ & $5.2{\times}10^{-1}$ & 66 & 10 & 140 & 150 \\
C & 31 & $10^{6}$ & BoomerAMG & 11 & $3.3{\times}10^{-9}$ & $1.2{\times}10^{-4}$ & 11 & 51 & 66 & 118 \\
C & 31 & $10^{6}$ & M2 $w{=}\frac1{32}$ & 12 & $1.7{\times}10^{-7}$ & $3.7{\times}10^{-3}$ & 14 & 14 & 34 & 48 \\
C & 31 & $10^{6}$ & M2 $w{=}\frac1{64}$ & 12 & $1.7{\times}10^{-7}$ & $3.7{\times}10^{-3}$ & 14 & 14 & 36 & 49 \\
C & 31 & $10^{8}$ & 传统 LGF & 500$^\times$ & $3.3{\times}10^{-1}$ & $9.1{\times}10^{-1}$ & -- & 0 & 688 & 688 \\
C & 31 & $10^{8}$ & GMG & 57 & $1.0{\times}10^{-7}$ & $4.6{\times}10^{-1}$ & 64 & 11 & 126 & 137 \\
C & 31 & $10^{8}$ & BoomerAMG & 11 & $1.0{\times}10^{-9}$ & $1.5{\times}10^{-4}$ & 10 & 59 & 75 & 134 \\
C & 31 & $10^{8}$ & M2 $w{=}\frac1{32}$ & 12 & $2.7{\times}10^{-7}$ & $8.4{\times}10^{-3}$ & 14 & 15 & 34 & 48 \\
C & 31 & $10^{8}$ & M2 $w{=}\frac1{64}$ & 12 & $2.7{\times}10^{-7}$ & $8.4{\times}10^{-3}$ & 14 & 20 & 47 & 67 \\
C & 63 & $10^{2}$ & 传统 LGF & 500$^\times$ & $1.6{\times}10^{-2}$ & $8.3{\times}10^{-1}$ & -- & 1 & 6175 & 6177 \\
C & 63 & $10^{2}$ & GMG & 89 & $1.4{\times}10^{-9}$ & $5.5{\times}10^{-1}$ & 86 & 82 & 1639 & 1721 \\
C & 63 & $10^{2}$ & BoomerAMG & 13 & $3.1{\times}10^{-10}$ & $1.8{\times}10^{-4}$ & 11 & 448 & 609 & 1057 \\
C & 63 & $10^{2}$ & M2 $w{=}\frac1{32}$ & 11 & $1.3{\times}10^{-8}$ & $1.2{\times}10^{-4}$ & 12 & 240 & 275 & 514 \\
C & 63 & $10^{2}$ & M2 $w{=}\frac1{64}$ & 16 & $1.4{\times}10^{-9}$ & $4.2{\times}10^{-4}$ & 15 & 129 & 346 & 475 \\
C & 63 & $10^{4}$ & 传统 LGF & 500$^\times$ & $2.2{\times}10^{-1}$ & $9.0{\times}10^{-1}$ & -- & 1 & 5015 & 5016 \\
C & 63 & $10^{4}$ & GMG & 88 & $2.9{\times}10^{-9}$ & $5.2{\times}10^{-1}$ & 86 & 88 & 1891 & 1979 \\
C & 63 & $10^{4}$ & BoomerAMG & 13 & $2.9{\times}10^{-9}$ & $6.7{\times}10^{-4}$ & 13 & 505 & 672 & 1176 \\
C & 63 & $10^{4}$ & M2 $w{=}\frac1{32}$ & 11 & $4.8{\times}10^{-8}$ & $4.7{\times}10^{-4}$ & 13 & 241 & 286 & 527 \\
C & 63 & $10^{4}$ & M2 $w{=}\frac1{64}$ & 15 & $1.7{\times}10^{-8}$ & $1.3{\times}10^{-3}$ & 16 & 125 & 312 & 437 \\
C & 63 & $10^{6}$ & 传统 LGF & 500$^\times$ & $3.3{\times}10^{-1}$ & $9.2{\times}10^{-1}$ & -- & 1 & 5179 & 5179 \\
C & 63 & $10^{6}$ & GMG & 78 & $1.5{\times}10^{-7}$ & $4.9{\times}10^{-1}$ & 85 & 89 & 1576 & 1665 \\
C & 63 & $10^{6}$ & BoomerAMG & 14 & $1.2{\times}10^{-9}$ & $6.3{\times}10^{-4}$ & 14 & 568 & 838 & 1406 \\
C & 63 & $10^{6}$ & M2 $w{=}\frac1{32}$ & 11 & $1.1{\times}10^{-7}$ & $1.4{\times}10^{-3}$ & 13 & 237 & 268 & 505 \\
C & 63 & $10^{6}$ & M2 $w{=}\frac1{64}$ & 15 & $4.8{\times}10^{-8}$ & $3.3{\times}10^{-3}$ & 17 & 125 & 332 & 457 \\
C & 63 & $10^{8}$ & 传统 LGF & 500$^\times$ & $4.1{\times}10^{-1}$ & $9.3{\times}10^{-1}$ & -- & 1 & 6245 & 6245 \\
C & 63 & $10^{8}$ & GMG & 75 & $2.2{\times}10^{-7}$ & $4.6{\times}10^{-1}$ & 84 & 88 & 1403 & 1491 \\
C & 63 & $10^{8}$ & BoomerAMG & 12 & $1.5{\times}10^{-8}$ & $8.1{\times}10^{-4}$ & 13 & 608 & 757 & 1365 \\
C & 63 & $10^{8}$ & M2 $w{=}\frac1{32}$ & 11 & $4.5{\times}10^{-7}$ & $2.7{\times}10^{-3}$ & 14 & 236 & 267 & 503 \\
C & 63 & $10^{8}$ & M2 $w{=}\frac1{64}$ & 15 & $3.1{\times}10^{-7}$ & $7.6{\times}10^{-3}$ & 18 & 125 & 304 & 429 \\
C & 127 & $10^{2}$ & 传统 LGF & 500$^\times$ & $2.0{\times}10^{-2}$ & $8.5{\times}10^{-1}$ & -- & 7 & 46400 & 46406 \\
C & 127 & $10^{2}$ & GMG & 103 & $3.1{\times}10^{-9}$ & $5.3{\times}10^{-1}$ & 99 & 795 & 14998 & 15793 \\
C & 127 & $10^{2}$ & BoomerAMG & 15 & $3.9{\times}10^{-10}$ & $7.2{\times}10^{-4}$ & 13 & 4135 & 5925 & 10059 \\
C & 127 & $10^{2}$ & M2 $w{=}\frac1{32}$ & 10 & $8.4{\times}10^{-9}$ & $6.1{\times}10^{-5}$ & 10 & 17148 & 3259 & 20407 \\
C & 127 & $10^{2}$ & M2 $w{=}\frac1{64}$ & 14 & $3.0{\times}10^{-9}$ & $2.5{\times}10^{-4}$ & 14 & 6043 & 3255 & 9298 \\
C & 127 & $10^{4}$ & 传统 LGF & 500$^\times$ & $2.5{\times}10^{-1}$ & $9.1{\times}10^{-1}$ & -- & 5 & 41668 & 41673 \\
C & 127 & $10^{4}$ & GMG & 103 & $5.0{\times}10^{-9}$ & $4.9{\times}10^{-1}$ & 100 & 641 & 14581 & 15222 \\
C & 127 & $10^{4}$ & BoomerAMG & 15 & $3.6{\times}10^{-10}$ & $9.1{\times}10^{-4}$ & 13 & 3467 & 5854 & 9320 \\
C & 127 & $10^{4}$ & M2 $w{=}\frac1{32}$ & 10 & $4.3{\times}10^{-8}$ & $2.4{\times}10^{-4}$ & 11 & 17309 & 3486 & 20795 \\
C & 127 & $10^{4}$ & M2 $w{=}\frac1{64}$ & 14 & $5.8{\times}10^{-9}$ & $8.7{\times}10^{-4}$ & 14 & 5946 & 3280 & 9226 \\
C & 127 & $10^{6}$ & 传统 LGF & 500$^\times$ & $3.6{\times}10^{-1}$ & $9.3{\times}10^{-1}$ & -- & 5 & 47432 & 47437 \\
C & 127 & $10^{6}$ & GMG & 101 & $2.8{\times}10^{-9}$ & $4.5{\times}10^{-1}$ & 100 & 773 & 14359 & 15133 \\
C & 127 & $10^{6}$ & BoomerAMG & 14 & $2.2{\times}10^{-9}$ & $1.0{\times}10^{-3}$ & 14 & 3750 & 5810 & 9560 \\
C & 127 & $10^{6}$ & M2 $w{=}\frac1{32}$ & 10 & $1.8{\times}10^{-7}$ & $6.4{\times}10^{-4}$ & 12 & 17751 & 3610 & 21361 \\
C & 127 & $10^{6}$ & M2 $w{=}\frac1{64}$ & 14 & $3.3{\times}10^{-8}$ & $2.0{\times}10^{-3}$ & 15 & 6283 & 3575 & 9857 \\
C & 127 & $10^{8}$ & 传统 LGF & 500$^\times$ & $4.4{\times}10^{-1}$ & $9.4{\times}10^{-1}$ & -- & 5 & 44078 & 44082 \\
C & 127 & $10^{8}$ & GMG & 89 & $2.9{\times}10^{-7}$ & $4.2{\times}10^{-1}$ & 99 & 638 & 12865 & 13502 \\
C & 127 & $10^{8}$ & BoomerAMG & 14 & $2.4{\times}10^{-9}$ & $2.3{\times}10^{-3}$ & 13 & 4059 & 6028 & 10087 \\
C & 127 & $10^{8}$ & M2 $w{=}\frac1{32}$ & 10 & $2.5{\times}10^{-7}$ & $1.2{\times}10^{-3}$ & 13 & 17589 & 3619 & 21208 \\
C & 127 & $10^{8}$ & M2 $w{=}\frac1{64}$ & 13 & $1.2{\times}10^{-7}$ & $3.9{\times}10^{-3}$ & 16 & 6042 & 3208 & 9249 \\
D & 31 & $10^{2}$ & 传统 LGF & 500$^\times$ & $3.1{\times}10^{-2}$ & $9.5{\times}10^{-1}$ & -- & 0 & 714 & 714 \\
D & 31 & $10^{2}$ & GMG & 88 & $2.0{\times}10^{-9}$ & $3.2{\times}10^{-1}$ & 86 & 10 & 197 & 207 \\
D & 31 & $10^{2}$ & BoomerAMG & 13 & $1.8{\times}10^{-10}$ & $4.5{\times}10^{-4}$ & 12 & 55 & 75 & 130 \\
D & 31 & $10^{2}$ & M2 $w{=}\frac1{32}$ & 13 & $2.8{\times}10^{-9}$ & $2.3{\times}10^{-4}$ & 13 & 10 & 37 & 47 \\
D & 31 & $10^{2}$ & M2 $w{=}\frac1{64}$ & 13 & $2.8{\times}10^{-9}$ & $2.3{\times}10^{-4}$ & 13 & 10 & 35 & 45 \\
D & 31 & $10^{4}$ & 传统 LGF & 500$^\times$ & $3.0{\times}10^{-1}$ & $9.8{\times}10^{-1}$ & -- & 0 & 733 & 733 \\
D & 31 & $10^{4}$ & GMG & 85 & $7.8{\times}10^{-9}$ & $2.9{\times}10^{-1}$ & 84 & 10 & 185 & 195 \\
D & 31 & $10^{4}$ & BoomerAMG & 14 & $1.6{\times}10^{-9}$ & $7.6{\times}10^{-4}$ & 13 & 59 & 87 & 147 \\
D & 31 & $10^{4}$ & M2 $w{=}\frac1{32}$ & 12 & $2.8{\times}10^{-8}$ & $3.6{\times}10^{-4}$ & 13 & 10 & 31 & 41 \\
D & 31 & $10^{4}$ & M2 $w{=}\frac1{64}$ & 12 & $2.8{\times}10^{-8}$ & $3.6{\times}10^{-4}$ & 13 & 10 & 29 & 39 \\
D & 31 & $10^{6}$ & 传统 LGF & 500$^\times$ & $4.4{\times}10^{-1}$ & $9.9{\times}10^{-1}$ & -- & 0 & 727 & 727 \\
D & 31 & $10^{6}$ & GMG & 74 & $1.1{\times}10^{-7}$ & $2.3{\times}10^{-1}$ & 80 & 10 & 159 & 168 \\
D & 31 & $10^{6}$ & BoomerAMG & 13 & $5.7{\times}10^{-9}$ & $1.2{\times}10^{-3}$ & 13 & 62 & 85 & 147 \\
D & 31 & $10^{6}$ & M2 $w{=}\frac1{32}$ & 12 & $3.9{\times}10^{-8}$ & $4.7{\times}10^{-4}$ & 14 & 10 & 30 & 40 \\
D & 31 & $10^{6}$ & M2 $w{=}\frac1{64}$ & 12 & $3.9{\times}10^{-8}$ & $4.7{\times}10^{-4}$ & 14 & 10 & 32 & 42 \\
D & 31 & $10^{8}$ & 传统 LGF & 500$^\times$ & $5.4{\times}10^{-1}$ & $9.9{\times}10^{-1}$ & -- & 0 & 710 & 710 \\
D & 31 & $10^{8}$ & GMG & 70 & $3.0{\times}10^{-7}$ & $1.6{\times}10^{-1}$ & 81 & 11 & 154 & 164 \\
D & 31 & $10^{8}$ & BoomerAMG & 14 & $6.6{\times}10^{-9}$ & $2.6{\times}10^{-3}$ & 14 & 62 & 93 & 155 \\
D & 31 & $10^{8}$ & M2 $w{=}\frac1{32}$ & 12 & $7.0{\times}10^{-8}$ & $7.7{\times}10^{-4}$ & 15 & 10 & 31 & 41 \\
D & 31 & $10^{8}$ & M2 $w{=}\frac1{64}$ & 12 & $7.0{\times}10^{-8}$ & $7.7{\times}10^{-4}$ & 15 & 10 & 32 & 42 \\
D & 63 & $10^{2}$ & 传统 LGF & 500$^\times$ & $6.0{\times}10^{-2}$ & $9.6{\times}10^{-1}$ & -- & 1 & 6486 & 6486 \\
D & 63 & $10^{2}$ & GMG & 107 & $3.6{\times}10^{-9}$ & $2.9{\times}10^{-1}$ & 104 & 99 & 1959 & 2058 \\
D & 63 & $10^{2}$ & BoomerAMG & 15 & $8.0{\times}10^{-11}$ & $8.5{\times}10^{-4}$ & 13 & 466 & 738 & 1204 \\
D & 63 & $10^{2}$ & M2 $w{=}\frac1{32}$ & 12 & $3.5{\times}10^{-9}$ & $9.6{\times}10^{-5}$ & 12 & 123 & 261 & 384 \\
D & 63 & $10^{2}$ & M2 $w{=}\frac1{64}$ & 17 & $5.0{\times}10^{-9}$ & $1.4{\times}10^{-3}$ & 17 & 71 & 346 & 417 \\
D & 63 & $10^{4}$ & 传统 LGF & 500$^\times$ & $3.9{\times}10^{-1}$ & $9.9{\times}10^{-1}$ & -- & 1 & 5440 & 5440 \\
D & 63 & $10^{4}$ & GMG & 105 & $1.4{\times}10^{-8}$ & $2.7{\times}10^{-1}$ & 109 & 90 & 1985 & 2075 \\
D & 63 & $10^{4}$ & BoomerAMG & 15 & $4.9{\times}10^{-10}$ & $9.3{\times}10^{-4}$ & 13 & 497 & 740 & 1237 \\
D & 63 & $10^{4}$ & M2 $w{=}\frac1{32}$ & 12 & $3.8{\times}10^{-9}$ & $1.5{\times}10^{-4}$ & 12 & 118 & 262 & 380 \\
D & 63 & $10^{4}$ & M2 $w{=}\frac1{64}$ & 17 & $7.9{\times}10^{-9}$ & $1.5{\times}10^{-3}$ & 17 & 67 & 307 & 375 \\
D & 63 & $10^{6}$ & 传统 LGF & 500$^\times$ & $5.4{\times}10^{-1}$ & $9.9{\times}10^{-1}$ & -- & 1 & 5318 & 5319 \\
D & 63 & $10^{6}$ & GMG & 100 & $6.0{\times}10^{-8}$ & $2.4{\times}10^{-1}$ & 106 & 94 & 1907 & 2001 \\
D & 63 & $10^{6}$ & BoomerAMG & 16 & $3.9{\times}10^{-9}$ & $2.6{\times}10^{-3}$ & 16 & 645 & 1039 & 1684 \\
D & 63 & $10^{6}$ & M2 $w{=}\frac1{32}$ & 12 & $5.7{\times}10^{-9}$ & $2.8{\times}10^{-4}$ & 12 & 111 & 239 & 350 \\
D & 63 & $10^{6}$ & M2 $w{=}\frac1{64}$ & 17 & $2.2{\times}10^{-8}$ & $2.2{\times}10^{-3}$ & 18 & 71 & 335 & 406 \\
D & 63 & $10^{8}$ & 传统 LGF & 500$^\times$ & $6.3{\times}10^{-1}$ & $9.9{\times}10^{-1}$ & -- & 1 & 6292 & 6293 \\
D & 63 & $10^{8}$ & GMG & 102 & $3.1{\times}10^{-8}$ & $2.0{\times}10^{-1}$ & 110 & 80 & 1846 & 1926 \\
D & 63 & $10^{8}$ & BoomerAMG & 15 & $1.3{\times}10^{-8}$ & $2.5{\times}10^{-3}$ & 16 & 625 & 1033 & 1659 \\
D & 63 & $10^{8}$ & M2 $w{=}\frac1{32}$ & 12 & $1.2{\times}10^{-8}$ & $4.1{\times}10^{-4}$ & 13 & 118 & 251 & 369 \\
D & 63 & $10^{8}$ & M2 $w{=}\frac1{64}$ & 17 & $7.9{\times}10^{-8}$ & $3.2{\times}10^{-3}$ & 19 & 70 & 306 & 376 \\
D & 127 & $10^{2}$ & 传统 LGF & 500$^\times$ & $6.7{\times}10^{-2}$ & $9.6{\times}10^{-1}$ & -- & 7 & 42085 & 42092 \\
D & 127 & $10^{2}$ & GMG & 135 & $1.4{\times}10^{-9}$ & $2.6{\times}10^{-1}$ & 129 & 640 & 19147 & 19787 \\
D & 127 & $10^{2}$ & BoomerAMG & 17 & $1.8{\times}10^{-10}$ & $2.0{\times}10^{-3}$ & 15 & 3862 & 7090 & 10953 \\
D & 127 & $10^{2}$ & M2 $w{=}\frac1{32}$ & 11 & $7.7{\times}10^{-9}$ & $5.9{\times}10^{-5}$ & 11 & 6552 & 2598 & 9150 \\
D & 127 & $10^{2}$ & M2 $w{=}\frac1{64}$ & 16 & $6.3{\times}10^{-9}$ & $6.7{\times}10^{-4}$ & 16 & 2693 & 2869 & 5562 \\
D & 127 & $10^{4}$ & 传统 LGF & 500$^\times$ & $4.1{\times}10^{-1}$ & $9.9{\times}10^{-1}$ & -- & 5 & 44985 & 44990 \\
D & 127 & $10^{4}$ & GMG & 133 & $1.9{\times}10^{-9}$ & $2.5{\times}10^{-1}$ & 129 & 829 & 20177 & 21006 \\
D & 127 & $10^{4}$ & BoomerAMG & 17 & $3.8{\times}10^{-10}$ & $4.4{\times}10^{-3}$ & 15 & 4073 & 7375 & 11448 \\
D & 127 & $10^{4}$ & M2 $w{=}\frac1{32}$ & 11 & $1.1{\times}10^{-8}$ & $9.9{\times}10^{-5}$ & 12 & 6863 & 2727 & 9590 \\
D & 127 & $10^{4}$ & M2 $w{=}\frac1{64}$ & 16 & $8.6{\times}10^{-9}$ & $7.9{\times}10^{-4}$ & 16 & 2760 & 2985 & 5745 \\
D & 127 & $10^{6}$ & 传统 LGF & 500$^\times$ & $5.6{\times}10^{-1}$ & $9.9{\times}10^{-1}$ & -- & 4 & 44751 & 44755 \\
D & 127 & $10^{6}$ & GMG & 128 & $3.4{\times}10^{-9}$ & $2.2{\times}10^{-1}$ & 126 & 667 & 18766 & 19433 \\
D & 127 & $10^{6}$ & BoomerAMG & 16 & $3.3{\times}10^{-9}$ & $3.0{\times}10^{-3}$ & 15 & 4153 & 7380 & 11532 \\
D & 127 & $10^{6}$ & M2 $w{=}\frac1{32}$ & 11 & $1.7{\times}10^{-8}$ & $2.0{\times}10^{-4}$ & 12 & 6695 & 2648 & 9344 \\
D & 127 & $10^{6}$ & M2 $w{=}\frac1{64}$ & 16 & $3.4{\times}10^{-8}$ & $1.5{\times}10^{-3}$ & 17 & 2683 & 2848 & 5531 \\
D & 127 & $10^{8}$ & 传统 LGF & 500$^\times$ & $6.6{\times}10^{-1}$ & $9.9{\times}10^{-1}$ & -- & 5 & 42997 & 43002 \\
D & 127 & $10^{8}$ & GMG & 123 & $9.9{\times}10^{-9}$ & $1.8{\times}10^{-1}$ & 123 & 746 & 18980 & 19726 \\
D & 127 & $10^{8}$ & BoomerAMG & 15 & $4.4{\times}10^{-9}$ & $1.5{\times}10^{-3}$ & 15 & 4657 & 6999 & 11656 \\
D & 127 & $10^{8}$ & M2 $w{=}\frac1{32}$ & 11 & $3.2{\times}10^{-8}$ & $2.9{\times}10^{-4}$ & 12 & 6868 & 2932 & 9801 \\
D & 127 & $10^{8}$ & M2 $w{=}\frac1{64}$ & 16 & $1.7{\times}10^{-7}$ & $2.1{\times}10^{-3}$ & 18 & 3032 & 2981 & 6013 \\
\bottomrule
\end{longtable}
\endgroup}{}

\section{复现}
代码在 \texttt{Nonlinear Poisson/FD\_LGF/}：\texttt{fdlgf\_m2.py}（\Mtwo{}）、\texttt{baselines/}（LGF、GMG、BoomerAMG）、
\texttt{benchmark.py}（二维）、\texttt{three\_d/}（三维全部代码）、\texttt{report/make\_tex.py}（本报告的表格）。
BoomerAMG 需要 \texttt{/opt/anaconda3/envs/hypre\_env}；测试：\texttt{test\_method2.py}、\texttt{test\_baselines.py}、\texttt{three\_d/test\_3d.py}。

\begin{thebibliography}{99}\small
\bibitem{xia} Q.~Xia, Unfitted boundary-algebraic preconditioning for exact-Jacobian Newton--Krylov solution of quasilinear elliptic problems in complex geometry, preprint (main.pdf).
\bibitem{alcouffe} R.~E.~Alcouffe, A.~Brandt, J.~E.~Dendy, J.~W.~Painter, The multi-grid method for the diffusion equation with strongly discontinuous coefficients, SIAM J. Sci. Stat. Comput. 2 (1981) 430--454.
\bibitem{agks} B.~Aksoylu, I.~G.~Graham, H.~Klie, R.~Scheichl, Towards a rigorously justified algebraic preconditioner for high-contrast diffusion problems, Comput. Visual. Sci. 11 (2008) 319--331.
\bibitem{briggs} W.~L.~Briggs, V.~E.~Henson, S.~F.~McCormick, A Multigrid Tutorial, 2nd ed., SIAM, 2000.
\bibitem{buzbee} B.~L.~Buzbee, F.~W.~Dorr, J.~A.~George, G.~H.~Golub, The direct solution of the discrete Poisson equation on irregular regions, SIAM J. Numer. Anal. 8 (1971) 722--736.
\bibitem{concus} P.~Concus, G.~H.~Golub, Use of fast direct methods for the efficient numerical solution of nonseparable elliptic equations, SIAM J. Numer. Anal. 10 (1973) 1103--1120.
\bibitem{gls} I.~G.~Graham, P.~O.~Lechner, R.~Scheichl, Domain decomposition for multiscale PDEs, Numer. Math. 106 (2007) 589--626.
\bibitem{henson} V.~E.~Henson, U.~M.~Yang, BoomerAMG: a parallel algebraic multigrid solver and preconditioner, Appl. Numer. Math. 41 (2002) 155--177.
\bibitem{lrt} R.~E.~Lynch, J.~R.~Rice, D.~H.~Thomas, Direct solution of partial difference equations by tensor product methods, Numer. Math. 6 (1964) 185--199.
\bibitem{mandel} J.~Mandel, Balancing domain decomposition, Comm. Numer. Methods Eng. 9 (1993) 233--241.
\bibitem{maury} B.~Maury, A fat boundary method for the Poisson problem in a domain with holes, J. Sci. Comput. 16 (2001) 319--339.
\bibitem{ms} H.~Moulinec, P.~Suquet, A numerical method for computing the overall response of nonlinear composites with complex microstructure, Comput. Methods Appl. Mech. Eng. 157 (1998) 69--94.
\bibitem{pw} W.~Proskurowski, O.~Widlund, On the numerical solution of Helmholtz's equation by the capacitance matrix method, Math. Comp. 30 (1976) 433--468.
\bibitem{tang} J.~M.~Tang, R.~Nabben, C.~Vuik, Y.~A.~Erlangga, Comparison of two-level preconditioners derived from deflation, domain decomposition and multigrid methods, J. Sci. Comput. 39 (2009) 340--370.
\bibitem{trottenberg} U.~Trottenberg, C.~W.~Oosterlee, A.~Sch\"uller, Multigrid, Academic Press, 2001.
\bibitem{xu} J.~Xu, Iterative methods by space decomposition and subspace correction, SIAM Rev. 34 (1992) 581--613.
\bibitem{zeman} J.~Zeman, J.~Vond\v{r}ejc, J.~Nov\'ak, I.~Marek, Accelerating a FFT-based solver for numerical homogenization of periodic media by conjugate gradients, J. Comput. Phys. 229 (2010) 8065--8071.
\end{thebibliography}
\end{document}
